 % Nombres en castellano (sin depender de babel)

% --- Caja de "Ideas clave" ---

% --- Teoremas en castellano, numerados por capítulo ---
\theoremstyle{plain}
\newtheorem{teorema}[theorem]{Teorema}
\newtheorem{proposicion}[theorem]{Proposición}
\newtheorem{corolario}[theorem]{Corolario}
\theoremstyle{definition}
\newtheorem{definicion}[theorem]{Definición}
\newtheorem{ejemplo}[theorem]{Ejemplo}
\theoremstyle{remark}
\newtheorem{observacion}[theorem]{Observación}

\title{Análisis Matricial --- Capítulo 5: Ortogonalidad}
\author{Sobre la base de Forzani, Joerin, Lovatto (2023/2025)\\
y \emph{Linear Algebra and Learning from Data}, G.\ Strang (Sección I.5)}
\date{2026}
 

 
\chapter{Ortogonalidad}\label{cap:ortogonalidad}

La palabra \textbf{ortogonal} significa \emph{perpendicular}, pero en álgebra
lineal su alcance es mucho mayor que el ángulo entre dos flechas. Lo que hace
este capítulo es seguir esa idea mientras crece, desde el caso más simple
hasta la factorización que la resume.

Se empieza por los \textbf{vectores}: dos vectores de $\R^n$ son ortogonales
cuando su producto escalar $\bx^T\by$ es cero, y ese único test da, casi
gratis, la desigualdad de Cauchy-Schwarz, la definición de ángulo y el teorema
de Pitágoras. De ahí se pasa a las \textbf{bases}: una base en la que todos los
vectores son perpendiculares entre sí y de longitud uno --- una base
\emph{ortonormal} --- vuelve triviales cuentas que en general son un sistema de
ecuaciones, porque cada coeficiente se obtiene con un producto escalar y sin
mirar a los demás.

El salto siguiente es de los vectores a los \textbf{subespacios}: dos
subespacios son ortogonales cuando cada vector de uno es perpendicular a cada
vector del otro. El ejemplo central ya lo teníamos sin saberlo, y es el que
completa el cuadro de los cuatro subespacios fundamentales del
Capítulo~\ref{cap:algebra-lineal}: el espacio fila de $A$ es ortogonal a su
espacio nulo, y el espacio columna lo es al espacio nulo de $A^T$.

Con subespacios ortogonales aparece la operación que domina el resto del
apunte: \textbf{proyectar}. Dado un subespacio y un vector de afuera, la
proyección es el punto del subespacio más cercano al vector, y resulta ser el
único para el cual el error es perpendicular al subespacio. Esa operación se
escribe como una matriz, y cuando el subespacio viene dado por una base
ortonormal escrita en las columnas de $Q$ esa matriz es simplemente $P=QQ^T$.

Las matrices $Q$ con columnas ortonormales, que cumplen $Q^TQ=I$, son las
protagonistas de la segunda mitad. Cuando además son cuadradas se llaman
\textbf{ortogonales} y satisfacen $Q^{-1}=Q^T$: multiplicar por ellas no
cambia longitudes ni ángulos, que es la razón por la que el cálculo numérico
las prefiere a cualquier otra. Veremos una colección de ejemplos --- rotaciones,
reflexiones, permutaciones, reflexiones de Householder, rotaciones de Givens,
matrices de Hadamard --- que reaparecen constantemente.

Todo esto se junta al final en la \textbf{factorización $A=QR$}, que escribe
cualquier matriz de columnas independientes como una matriz de columnas
ortonormales por una triangular superior. La construiremos de dos maneras: con
el proceso de Gram-Schmidt, que es el camino conceptual, y con reflexiones de
Householder, que es el que se usa en la práctica. Los mínimos cuadrados, que
son la gran aplicación de las proyecciones y de $QR$, tienen capítulo propio
más adelante.

% ===================================================================
\section{Vectores ortogonales}

Las dos operaciones con las que vamos a trabajar ya están definidas: el
\textbf{producto escalar} y la \textbf{norma euclídea} de la
Sección~\ref{sec:producto-escalar} del Capítulo~\ref{cap:algebra-lineal},
\[
\bx^T\by=x_1y_1+x_2y_2+\cdots+x_ny_n ,
\qquad
\|\bx\|=\sqrt{\bx^T\bx}=\sqrt{x_1^2+\cdots+x_n^2} ,
\]
junto con la noción de vectores \textbf{ortogonales}, que son los que cumplen
$\bx^T\by=0$ y que anotamos $\bx\perp\by$. Un vector es \textbf{unitario} si
$\|\bx\|=1$, y el vector nulo es ortogonal a todos.

Allá dijimos que las propiedades quedaban para este capítulo, y de eso se
trata esta sección: no de definir nada nuevo, sino de probar lo que esas dos
operaciones cumplen. Dos propiedades son inmediatas desde la definición y
conviene tenerlas a mano. Del producto escalar:
\begin{enumerate}[label=\textup{(\arabic*)},itemsep=1pt,topsep=3pt]
\item $\bx^T\bx\ge0$, y $\bx^T\bx=0$ si y sólo si $\bx=\bi{0}$;
\item $\bx^T\by=\by^T\bx$ (es un número, luego igual a su traspuesta);
\item $(a\bx+\by)^T\bu=a\,\bx^T\bu+\by^T\bu$ para todo $a\in\R$: el producto
escalar es lineal en cada variable.
\end{enumerate}
Y de la norma, las dos primeras propiedades de las tres que la definen:
$\|\bx\|\ge0$ con igualdad sólo en $\bx=\bi{0}$, y $\|a\bx\|=|a|\,\|\bx\|$.

La tercera, la \textbf{desigualdad triangular}
$\|\bx+\by\|\le\|\bx\|+\|\by\|$, no es inmediata: necesita la desigualdad de
Cauchy-Schwarz, que probamos a continuación y que es además la que permite
definir el \emph{ángulo} entre dos vectores. Después de eso viene el teorema
de Pitágoras, y con esas tres herramientas queda armado todo lo que el resto
del capítulo necesita.

\begin{teorema}[Desigualdad de Cauchy-Schwarz]\label{teo:cauchy-schwarz}
Para todos $\bx,\by\in\R^n$,
\[
|\bx^T\by|\le\|\bx\|\,\|\by\| ,
\]
con igualdad si y solo si $\bx$ e $\by$ son linealmente dependientes.
\end{teorema}

\begin{proof}
Si $\bx=\bi{0}$ no hay nada que probar. Si no, miramos la función
cuadrática en $t$
\[
p(t)=\|t\bx+\by\|^2=(t\bx+\by)^T(t\bx+\by)
=t^2\,\|\bx\|^2+2t\,\bx^T\by+\|\by\|^2 .
\]
Como $p(t)\ge0$ para todo $t$ (¡es una norma al cuadrado!), la parábola no
puede tener dos raíces reales distintas, así que su discriminante es
$\le 0$:
\[
4(\bx^T\by)^2-4\,\|\bx\|^2\|\by\|^2\le 0 .
\]
Hay igualdad exactamente cuando $p$ tiene una raíz doble $t_0$, es decir
cuando $\|t_0\bx+\by\|=0$: cuando $\by=-t_0\bx$ es múltiplo de $\bx$.
\end{proof}

\begin{proposicion}[Desigualdad triangular]
$\|\bx+\by\|\le\|\bx\|+\|\by\|$ para todos $\bx,\by\in\R^n$.
\end{proposicion}

\begin{proof}
$\|\bx+\by\|^2=\|\bx\|^2+2\,\bx^T\by+\|\by\|^2
\le\|\bx\|^2+2\|\bx\|\|\by\|+\|\by\|^2=(\|\bx\|+\|\by\|)^2$,
usando Cauchy-Schwarz en el término del medio.
\end{proof}

Cauchy-Schwarz dice que el cociente $\bx^T\by/(\|\bx\|\|\by\|)$
siempre cae en $[-1,1]$ (para vectores no nulos): por lo tanto es el coseno
de un ángulo. \emph{Definimos} el ángulo $\theta$ entre $\bx$ e $\by$ por
\[
\cos\theta=\frac{\bx^T\by}{\|\bx\|\,\|\by\|}\ ,
\]
que coincide con el ángulo geométrico en el plano que generan. Ser
ortogonales es exactamente tener $\cos\theta=0$, es decir
$\theta=90^{\circ}$.

\begin{observacion}[Similitud coseno]
En ciencia de datos, ese coseno se usa constantemente como medida de
\emph{similitud} entre dos vectores (documentos representados por conteos
de palabras, \emph{embeddings}, perfiles de usuarios): la
\textbf{similitud coseno} compara direcciones ignorando las magnitudes.
Dos vectores con las mismas proporciones tienen similitud $1$ aunque uno
sea ``el doble de largo'' que el otro (ese es justo el caso de igualdad
de Cauchy-Schwarz). Los ortogonales tienen similitud $0$: no comparten
información direccional.
\end{observacion}

\begin{teorema}[Pitágoras]\label{teo:pitagoras}
Si $\bx\perp\by$ entonces
\[
\|\bx-\by\|^2=\|\bx\|^2+\|\by\|^2
\qquad\text{y también}\qquad
\|\bx+\by\|^2=\|\bx\|^2+\|\by\|^2 .
\]
\end{teorema}

\begin{proof}
Desarrollamos el cuadrado usando que el producto interno es lineal:
\[
\|\bx-\by\|^2=(\bx-\by)^T(\bx-\by)
=\bx^T\bx-\bx^T\by-\by^T\bx+\by^T\by
=\|\bx\|^2+\|\by\|^2-2\,\bx^T\by .
\]
Si $\bx^T\by=0$, los dos términos cruzados desaparecen y queda
Pitágoras. El mismo cálculo con $\bx+\by$ da el otro enunciado.
\end{proof}

\begin{ejemplo}
$\bx=(1,2,2)^T$ y $\by=(2,1,-2)^T$ cumplen
$\bx^T\by=2+2-4=0$. La hipotenusa es $\bx-\by=(-1,1,4)^T$ y
Pitágoras dice $18=9+9$. \hfill$\square$
\end{ejemplo}

\begin{observacion}[Ley de cosenos]
Sin suponer ortogonalidad, el cálculo de la demostración anterior da
\[
\|\bx-\by\|^2=\|\bx\|^2+\|\by\|^2-2\,\|\bx\|\,\|\by\|\cos\theta ,
\]
la ley de cosenos $c^2=a^2+b^2-2ab\cos\theta$. Los vectores ortogonales
tienen $\cos\theta=0$ y ese último término desaparece.
\end{observacion}


\subsection*{El caso complejo}\label{sec:producto-interno-complejo}

Para vectores complejos $\bx,\by\in\Cx^n$ el test cambia: hay que conjugar.
El producto interno usual se define como
\[
\overline{\bx}^T\by=\overline{x}_1y_1+\cdots+\overline{x}_ny_n ,
\qquad
\|\bx\|^2=\overline{\bx}^T\bx=|x_1|^2+\cdots+|x_n|^2 ,
\]
y $\bx\perp\by$ significa $\overline{\bx}^T\by=0$.

¿Por qué conjugar? Porque sin conjugar la norma deja de servir. Lo único que
le pedimos a $\|\bx\|^2$ es que sea un número \emph{real y no negativo}, que
valga cero sólo en el vector nulo; eso es lo que permite hablar de longitudes
y de distancias. Si copiáramos la fórmula real sin cambios, tendríamos
$\bx^T\bx=\sum_k x_k^2$, y los cuadrados de números complejos no tienen por qué
ser positivos ni siquiera reales. El ejemplo más chico lo muestra: para
$\bx=(1,i)^T$ es
\[
\bx^T\bx = 1^2+i^2 = 1-1 = 0 ,
\]
o sea un vector no nulo tendría longitud cero. En cambio, conjugando el
primer factor cada término se vuelve
$\overline{x}_k x_k=|x_k|^2\ge0$, y entonces
\[
\overline{\bx}^T\bx=|x_1|^2+\cdots+|x_n|^2\ \ge 0 ,
\]
que se anula sólo si todas las coordenadas son nulas. Para el mismo ejemplo,
$\overline{\bx}^T\bx=|1|^2+|i|^2=2$ y $\|\bx\|=\sqrt2$, que es lo razonable.

El precio es que el producto deja de ser simétrico: ahora
$\overline{\by}^T\bx=\overline{\overline{\bx}^T\by}$, es decir, intercambiar
los argumentos conjuga el resultado. Y la linealidad vale en la segunda
variable pero en la primera aparece un conjugado:
$\overline{(c\bx)}^T\by=\overline{c}\,\overline{\bx}^T\by$. Todo lo demás ---
Cauchy-Schwarz, Pitágoras, la desigualdad triangular, las proyecciones ---
sigue valiendo palabra por palabra, cambiando $\bx^T$ por $\overline{\bx}^T$.
La matriz de Fourier de los ejercicios es el ejemplo estrella de ortogonalidad
compleja.


% ===================================================================
\section{Bases ortogonales y ortonormales}

\begin{definicion}
Una base $\bv_1,\dots,\bv_r$ de un subespacio \textbf{de} $\R^n$ es una \textbf{base ortogonal}
si $\bv_i^T\bv_j=0$ para todo $i\ne j$. Si además todos los vectores son
unitarios ($\bv_i^T\bv_i=1$), es una \textbf{base ortonormal}. En ese
caso escribimos $\bq_1,\dots,\bq_r$:
\[
\bq_i^T\bq_j=\begin{cases}0 & i\ne j\\ 1 & i=j.\end{cases}
\]
\end{definicion}

De ortogonal a ortonormal se pasa dividiendo cada vector por su norma:
$\bq_i=\bv_i/\|\bv_i\|$. El ejemplo más simple es la \textbf{base canónica}
$\bi{e}_1,\dots,\bi{e}_n$ de $\R^n$, que es ortonormal: en $\R^3$,
\[
\bi{i}=\begin{bmatrix}1\\0\\0\end{bmatrix}\quad
\bi{j}=\begin{bmatrix}0\\1\\0\end{bmatrix}\quad
\bi{k}=\begin{bmatrix}0\\0\\1\end{bmatrix}.
\]

\begin{proposicion}\label{prop:ort-indep}
Vectores ortogonales no nulos son linealmente independientes.
\end{proposicion}

\begin{proof}
Supongamos $c_1\bv_1+\cdots+c_r\bv_r=\bi{0}$ con los $\bv_i$ ortogonales no
nulos. Multiplicando por $\bv_k^T$ a la izquierda, todos los productos
$\bv_k^T\bv_i$ con $i\ne k$ se anulan y queda
$c_k\,\bv_k^T\bv_k=0$. Como $\bv_k\ne\bi{0}$, es $\|\bv_k\|^2>0$ y por
lo tanto $c_k=0$, para cada $k$.
\end{proof}

?` Cu\'al es el an\'alogo en $\Cx^n$?

\subsection*{Coeficientes en una base ortonormal}

Aquí está la gran ventaja práctica de la ortonormalidad. Si
$\bq_1,\dots,\bq_n$ es una base ortonormal de $\R^n$, todo vector se escribe
\begin{equation}\label{eq:combinacion}
\bv=c_1\bq_1+c_2\bq_2+\cdots+c_n\bq_n .
\end{equation}
En una base cualquiera, encontrar los $c_i$ exige resolver un sistema lineal.
En una base ortonormal, \emph{cada coeficiente se obtiene por separado con un
producto interno}:

\begin{teorema}[Coeficientes en una base ortonormal]\label{teo:coeficientes}
\[
c_1=\bq_1^T\bv,\qquad c_2=\bq_2^T\bv,\qquad\dots\qquad
c_n=\bq_n^T\bv .
\]
\end{teorema}

\begin{proof}
\emph{Demostración vectorial.} Multiplicamos \eqref{eq:combinacion} por
$\bq_1^T$:
\[
\bq_1^T\bv=c_1\,\bq_1^T\bq_1+c_2\,\bq_1^T\bq_2+\cdots
+c_n\,\bq_1^T\bq_n=c_1 ,
\]
porque todos los términos se anulan salvo $c_1\bq_1^T\bq_1=c_1$. Lo
mismo con cada $\bq_k$.

\emph{Demostración matricial.} La ecuación \eqref{eq:combinacion} es
$\bv=Q\bi{c}$ con $Q=[\bq_1\ \cdots\ \bq_n]$. Multiplicando por $Q^T$:
\[
Q^T\bv=Q^T Q\,\bi{c}=\bi{c},
\]
que da todos los coeficientes $c_k=\bq_k^T\bv$ de una sola vez.
\end{proof}

\begin{observacion}
Esta es la aplicación clave de las bases ortonormales (por ejemplo, la base
de las series de Fourier): cada coeficiente se calcula de manera
independiente de los demás. Los términos $c_k\bq_k$ son las componentes de
$\bv$ a lo largo de los ejes ortogonales definidos por la base: son las
proyecciones de $\bv$ sobre esos ejes.
\end{observacion}

\

?` Cu\'al es el an\'alogo en $\Cx^n$?

\subsection*{Matrices de Hadamard}

¿Existen bases ortogonales de $\R^n$ formadas solo por $1$ y $-1$? Para
$n=2,4,8$ sí, y son las columnas de las \textbf{matrices de Hadamard}:
\[
H_2=\begin{bmatrix}1&1\\1&-1\end{bmatrix}\qquad
H_4=\begin{bmatrix}1&1&1&1\\1&-1&1&-1\\1&1&-1&-1\\1&-1&-1&1\end{bmatrix}
\qquad
H_8=\begin{bmatrix}H_4&H_4\\ H_4&-H_4\end{bmatrix}.
\]
¿Son matrices ortogonales? \emph{No todavía}: sus columnas tienen normas
$\sqrt2,\sqrt4,\sqrt8$. Dividiendo por esas normas obtenemos bases
ortonormales de $\R^2,\R^4,\R^8$, y la construcción por bloques
$H_{2n}=\left[\begin{smallmatrix}H_n&H_n\\H_n&-H_n\end{smallmatrix}\right]$
continúa la lista para $n=16,32,\dots$

\begin{observacion}[Conjetura de Hadamard]
Si $H_n$ existe (con $n>2$), entonces $n$ debe ser múltiplo de $4$ (ver
ejercicios). La \textbf{conjetura de Hadamard} propone la recíproca: que hay
una matriz de $\pm1$ con columnas ortogonales siempre que $4$ divide a $n$.
Sigue abierta; el tamaño más chico sin matriz de Hadamard conocida es
$n=668$.
\end{observacion}

% ===================================================================
\section{Subespacios ortogonales en $\R^n$ con el producto interno usual}\label{sec:cuatro-ort}

\begin{definicion}
Dos subespacios $\mathbf{V}$ y $\mathbf{W}$ de $\R^n$ son
\textbf{ortogonales} si todo vector de $\mathbf{V}$ es ortogonal a todo
vector de $\mathbf{W}$: $\bv^T\bw=0$ para todo
$\bv\in\mathbf{V},\ \bw\in\mathbf{W}$.
El \textbf{complemento ortogonal} de $\mathbf{V}$ es
\[
\mathbf{V}^{\perp}=\{\bw\in\R^n:\ \bw^T\bv=0
\text{ para todo }\bv\in\mathbf{V}\}.
\]
\end{definicion}

En $\R^3$, dos paredes perpendiculares \emph{no} son subespacios ortogonales:
comparten la recta de intersección, y un vector no nulo nunca es ortogonal a
sí mismo. Dos subespacios ortogonales solo se cortan en $\{\bi{0}\}$ (por
Pitágoras, o por la Proposición~\ref{prop:ort-indep}).

\begin{observacion}[Ortogonales no es lo mismo que complementos]
Dos subespacios ortogonales no son necesariamente complementos ortogonales
uno del otro. En $\R^3$, $\mathbf{U}=\hbox{span}\{(1,0,0)\}$ y
$\mathbf{W}=\hbox{span}\{(0,1,0)\}$ son ortogonales, pero el complemento
ortogonal de $\mathbf{U}$ es
$\mathbf{U}^{\perp}=\hbox{span}\{(0,1,0),(0,0,1)\}$, que contiene
\emph{estrictamente} a $\mathbf{W}$: falta una dimensión. Ser
complementos exige más que ser ortogonales: las dimensiones tienen que
sumar $n$.
\end{observacion}

\begin{proposicion}[Descomposición ortogonal]\label{prop:descomposicion}
Sea $\mathbf{W}$ un subespacio de $\R^n$. Entonces:
\begin{enumerate}[label=(\roman*),itemsep=1pt]
\item $\mathbf{W}\cap\mathbf{W}^{\perp}=\{\bi{0}\}$;
\item todo $\bv\in\R^n$ se escribe de manera \emph{única} como
$\bv=\bw+\bi{z}$ con $\bw\in\mathbf{W}$ y $\bi{z}\in\mathbf{W}^{\perp}$
(escribimos $\R^n=\mathbf{W}\oplus\mathbf{W}^{\perp}$);
\item $\dim\mathbf{W}^{\perp}=n-\dim\mathbf{W}$;
\item $(\mathbf{W}^{\perp})^{\perp}=\mathbf{W}$.
\end{enumerate}
\end{proposicion}

\begin{proof}
(i) Si $\bw$ está en ambos, es ortogonal a sí mismo: $\bw^T\bw=0$ y
$\bw=\bi{0}$.

(ii) \emph{Existencia:} tomemos una base ortonormal $\bq_1,\dots,\bq_r$ de
$\mathbf{W}$ (existe por Gram-Schmidt, Sección~\ref{sec:gram-schmidt}, que
no usa nada de esta proposición). Sea
$\bw=\bq_1(\bq_1^T\bv)+\cdots+\bq_r(\bq_r^T\bv)\in\mathbf{W}$ y
$\bi{z}=\bv-\bw$. Para cada $k$,
$\bq_k^T\bi{z}=\bq_k^T\bv-\bq_k^T\bv=0$, y como todo
elemento de $\mathbf{W}$ es combinación de los $\bq_k$, resulta
$\bi{z}\perp\mathbf{W}$. \emph{Unicidad:} si
$\bw_1+\bi{z}_1=\bw_2+\bi{z}_2$, entonces
$\bw_1-\bw_2=\bi{z}_2-\bi{z}_1\in\mathbf{W}\cap\mathbf{W}^{\perp}
=\{\bi{0}\}$.

(iii) Juntando una base de $\mathbf{W}$ con una de $\mathbf{W}^{\perp}$ se
obtiene, por (ii), un conjunto generador de $\R^n$ que es independiente
(Proposición~\ref{prop:ort-indep} aplicada por bloques, o directamente por
(i)): una base. Las dimensiones suman $n$.

(iv) $\mathbf{W}\subseteq(\mathbf{W}^{\perp})^{\perp}$ es inmediato de la
definición. Recíprocamente, si $\bv\in(\mathbf{W}^{\perp})^{\perp}$,
escribimos $\bv=\bw+\bi{z}$ como en (ii); entonces
$\bi{z}^T\bi{z}=\bi{z}^T(\bv-\bw)
=\bi{z}^T\bv-\bi{z}^T\bw=0-0=0$
(el primer término se anula porque $\bi{z}\in\mathbf{W}^{\perp}$ y $\bv$ es
ortogonal a todo $\mathbf{W}^{\perp}$). Luego $\bi{z}=\bi{0}$ y
$\bv=\bw\in\mathbf{W}$.
\end{proof}

El ejemplo que organiza todo el curso viene de $A\bx=\bi{0}$:

\begin{teorema}[Espacio fila $\perp$ espacio nulo]\label{teo:fila-nulo}
Para toda matriz $A$ de $m\times n$:
\[
C(A^T)\perp N(A)
\qquad\text{y}\qquad
C(A)\perp N(A^T).
\]
Más aún, $N(A)=C(A^T)^{\perp}$ y $N(A^T)=C(A)^{\perp}$.
\end{teorema}

\begin{proof}
La definición del espacio nulo ya contiene la ortogonalidad. Que
$A\bx=\bi{0}$ significa que \emph{cada fila} de $A$ multiplica a $\bx$ y da
cero:
\[
A\bx=
\begin{bmatrix}\text{fila }1\text{ de }A\\ \vdots\\ \text{fila }m\text{ de }A
\end{bmatrix}
\begin{bmatrix}\ \\ \bx\\ \ \end{bmatrix}
=\begin{bmatrix}0\\\vdots\\0\end{bmatrix}.
\]
Entonces $\bx$ es ortogonal a cada fila, y por linealidad a cada combinación
de filas: a todo el espacio fila $C(A^T)$. Recíprocamente, si $\bx$ es
ortogonal a todas las filas, entonces $A\bx=\bi{0}$; por eso $N(A)$ es
\emph{exactamente} el complemento ortogonal del espacio fila.

Aplicando lo mismo a $A^T\by=\bi{0}$: las columnas de $A$ (filas de
$A^T$) son todas ortogonales a $\by$, y también sus combinaciones. El
espacio columna de $A$ es ortogonal al espacio nulo de $A^T$.
\end{proof}

\begin{observacion}[Las dimensiones encajan]
Las dimensiones $r$ y $n-r$ suman $n$: una base del espacio fila ($r$
vectores) junto con una base del espacio nulo ($n-r$ vectores) forma una base
de todo $\R^n$. En consecuencia, \textbf{todo vector $\bv\in\R^n$ se
descompone de manera única} como
\[
\bv=\bv_{\text{fila}}+\bv_{\text{nulo}},
\qquad \bv_{\text{fila}}\in C(A^T),\quad \bv_{\text{nulo}}\in N(A),
\]
y las dos componentes son ortogonales entre sí. Lo mismo pasa en $\R^m$ con
$C(A)$ y $N(A^T)$, de dimensiones $r$ y $m-r$.
\end{observacion}

\begin{figure}[ht]
\centering
\begin{tikzpicture}[scale=0.9,>=stealth]
% --- lado Rn ---
\begin{scope}[shift={(0,0)}]
\draw[thick,rotate=45] (-1.5,-1.05) rectangle (1.5,1.05);
\node at (0,0.55) {\small espacio};
\node at (0,0.05) {\small fila de $A$};
\node at (0,-0.5) {\small $\dim=r$};
\end{scope}
\begin{scope}[shift={(1.15,-3.2)}]
\draw[thick,rotate=-45] (-1.4,-0.95) rectangle (1.4,0.95);
\node at (0,0.45) {\small espacio};
\node at (0,-0.05) {\small nulo de $A$};
\node at (0,-0.55) {\small $\dim=n-r$};
\end{scope}
\node at (-1.9,-2.4) {$\R^n$};
\draw (0.72,-1.28) rectangle +(0.28,0.28); % angulo recto
\node at (1.45,-1.05) {\small $90^{\circ}$};
% --- lado Rm ---
\begin{scope}[shift={(8.3,0.2)}]
\draw[thick,rotate=45] (-1.5,-1.05) rectangle (1.5,1.05);
\node at (0,0.55) {\small espacio};
\node at (0,0.05) {\small columna de $A$};
\node at (0,-0.5) {\small $\dim=r$};
\end{scope}
\begin{scope}[shift={(9.45,-3.0)}]
\draw[thick,rotate=-45] (-1.4,-0.95) rectangle (1.4,0.95);
\node at (0,0.45) {\small espacio nulo};
\node at (0,-0.05) {\small de $A^T$};
\node at (0,-0.55) {\small $\dim=m-r$};
\end{scope}
\node at (11.3,-2.2) {$\R^m$};
\draw (9.02,-1.08) rectangle +(0.28,0.28);
\node at (8.35,-0.9) {\small $90^{\circ}$};
% --- flechas ---
\draw[->,thick] (2.4,0.4) -- node[above]{\small $A\bx_{\text{fila}}=\bi{b}$} (6.1,0.5);
\draw[->,thick] (3.3,-2.6) -- node[below]{\small $A\bx_{\text{nulo}}=\bi{0}$} (6.6,-1.4);
\end{tikzpicture}
\caption{Dos pares de subespacios ortogonales. Las dimensiones suman $n$ y
suman $m$. Esta es la \emph{gran figura} del álgebra lineal: dos subespacios
en $\R^n$ y dos en $\R^m$.}
\label{fig:granfigura}
\end{figure}

\begin{corolario}[Los cuatro subespacios, completo]\label{coro:cuatro}
Sea $A\in\R^{m\times n}$ de rango $r$. Entonces
\[
\R^n = C(A^T)\oplus N(A),
\qquad
\R^m = C(A)\oplus N(A^T),
\]
y ambas descomposiciones son ortogonales. En particular
\[
\dim C(A^T)=r,\quad \dim N(A)=n-r,\qquad
\dim C(A)=r,\quad \dim N(A^T)=m-r,
\]
y $\big(C(A^T)\big)^{\perp\perp}=C(A^T)$, $\big(N(A)\big)^{\perp}=C(A^T)$.
\end{corolario}

\begin{proof}
Por el Teorema~\ref{teo:fila-nulo}, $N(A)=C(A^T)^{\perp}$; la
Proposición~\ref{prop:descomposicion} aplicada a $\mathbf{W}=C(A^T)$ da la
suma directa ortogonal y $\dim N(A)=n-\dim C(A^T)=n-r$. Lo mismo con $A^T$
en lugar de $A$.
\end{proof}

\begin{observacion}[Dónde se calcula esta descomposición]
El Corolario~\ref{coro:cuatro} dice que todo $\bx\in\R^n$ se parte de manera
única en $\bx=\bx_{\text{fila}}+\bx_{\text{nulo}}$, pero no dice cómo
calcular las dos piezas. La respuesta está en la sección siguiente: si $P$ es
la proyección ortogonal sobre el espacio fila, entonces
$\bx_{\text{fila}}=P\bx$ y $\bx_{\text{nulo}}=(I-P)\bx$. Las proyecciones
son la herramienta \emph{computacional} que corresponde a esta figura.
\end{observacion}

\begin{observacion}[Dos consecuencias que usaremos todo el tiempo]
\label{obs:ATA}
De $N(A)=C(A^T)^{\perp}$ se deduce en dos líneas que
\[
N(A^TA)=N(A)
\qquad\text{y}\qquad
C(A^TA)=C(A^T),
\]
donde la primera igualdad ya apareció en el
Ejercicio de desarrollo~\ref{ent:rango} del Capítulo~\ref{cap:algebra-lineal}
y dice, en particular, que \emph{$A^TA$ es invertible si y sólo si las
columnas de $A$ son independientes}. Esa es la hipótesis que hará falta para
proyectar sobre $C(A)$ cuando la base no sea ortonormal.
\end{observacion}

\begin{observacion}[Anticipo: las bases especiales de la SVD]
Gram-Schmidt (Sección~\ref{sec:gram-schmidt}) produce bases ortonormales
$\bv_1,\dots,\bv_r$ del espacio fila y $\bu_1,\dots,\bu_r$ del espacio
columna. La descomposición en valores singulares (SVD), el teorema más
importante de la ciencia de datos, encuentra bases con una propiedad extra:
cada par queda conectado por $A$,
\[
A\bv_1=\sigma_1\bu_1,\quad A\bv_2=\sigma_2\bu_2,\quad\dots\quad
A\bv_r=\sigma_r\bu_r .
\]
En la Figura~\ref{fig:granfigura}: multiplicar por $A$ lleva una base
ortogonal de $\bv$'s (a la izquierda) a una base ortogonal de $\bu$'s (a la
derecha). Volveremos sobre esto.
\end{observacion}

% ===================================================================
\section{Proyecciones ortogonales en $\R^n$ con el producto interno usual}\label{sec:proyecciones}

\subsection*{Proyección sobre una recta}

Sea $\ba\ne\bi{0}$ y sea la recta $\gen\{\ba\}$. La \textbf{proyección
ortogonal} de $\bb$ sobre esa recta es el punto de la recta más cercano a
$\bb$. Todo punto de la recta es un múltiplo $\alpha\,\ba$ con
$\alpha\in\R$, así que buscamos el \emph{escalar} $\alpha$ para el cual el
\emph{error} $\bi{e}=\bb-\alpha\,\ba$ resulta ortogonal a $\ba$:
\[
\ba^T(\bb-\alpha\,\ba)=0
\qquad\Longrightarrow\qquad
\alpha=\frac{\ba^T\bb}{\ba^T\ba}
\qquad\Longrightarrow\qquad
\boxed{\ \bi{p}=\frac{\ba^T\bb}{\ba^T\ba}\,\ba\ }
\]
En forma matricial, $\bi{p}=P\bb$ con
$P=\dfrac{\ba\,\ba^T}{\ba^T\ba}$, una matriz de rango uno. Si el
vector es unitario, $\bq=\ba/\|\ba\|$, la fórmula es simplemente
$P=\bq\bq^T$ y $\bi{p}=\bq(\bq^T\bb)$: el coeficiente
$\bq^T\bb$ del Teorema~\ref{teo:coeficientes} multiplicado por $\bq$.

\subsection*{Proyección sobre un subespacio: $P=QQ^T$}

Sea $\mathbf{V}$ un subespacio de $\R^m$ con base ortonormal
$\bq_1,\dots,\bq_r$, y sea $Q=[\bq_1\ \cdots\ \bq_r]$ de $m\times r$, de modo
que $Q^T Q=I_r$. Con una base ortonormal, proyectar sobre $\mathbf{V}$
es proyectar sobre cada eje y sumar:
\[
\bi{p}=\bq_1(\bq_1^T\bb)+\cdots+\bq_r(\bq_r^T\bb)
=QQ^T\bb .
\]

\begin{teorema}\label{teo:proyector}
La matriz $P=QQ^T$ satisface
\[
P=P^T=P^2 ,
\]
y $P\bb$ es la proyección ortogonal de $\bb$ sobre $C(Q)$: el error
$\bi{e}=\bb-P\bb$ es ortogonal a $C(Q)$, y $P\bb$ es el punto de $C(Q)$ más
cercano a $\bb$. Recíprocamente, toda matriz con $P^2=P=P^T$ es la
proyección ortogonal sobre su espacio columna.
\end{teorema}

\begin{proof}
Simetría: $(QQ^T)^T=QQ^T$. Idempotencia: en el medio
eliminamos $Q^T Q=I$,
\[
P^2=(QQ^T)(QQ^T)=Q\,(Q^T Q)\,Q^T=QQ^T=P .
\]
El error es ortogonal al subespacio: para cada columna $\bq_k$,
\[
Q^T(\bb-QQ^T\bb)=Q^T\bb-(Q^T Q)Q^T\bb
=Q^T\bb-Q^T\bb=\bi{0}.
\]
Y estar más cerca es Pitágoras: si $\bv\in C(Q)$ es cualquier otro punto,
$\bb-\bv=(\bb-P\bb)+(P\bb-\bv)$ con los dos sumandos ortogonales
(el primero es el error, el segundo vive en $C(Q)$), así que
$\|\bb-\bv\|^2=\|\bb-P\bb\|^2+\|P\bb-\bv\|^2\ge\|\bb-P\bb\|^2$.

\emph{Recíproca.} Sea ahora $P$ cualquier matriz con $P^2=P=P^T$ y sea
$\mathbf{W}=C(P)$. Para todo $\bb$: primero, $P\bb\in\mathbf{W}$. Segundo,
el ``error'' $\bb-P\bb$ está en el espacio nulo de $P$, porque
$P(\bb-P\bb)=P\bb-P^2\bb=\bi{0}$; y por simetría
$N(P)=N(P^T)=C(P)^{\perp}$ (Teorema~\ref{teo:fila-nulo}). Es decir,
$P\bb\in\mathbf{W}$ y $\bb-P\bb\perp\mathbf{W}$: por la unicidad de la
descomposición ortogonal (Proposición~\ref{prop:descomposicion}), $P\bb$
es \emph{la} proyección ortogonal de $\bb$ sobre $\mathbf{W}$. Las dos
condiciones $P^2=P$ y $P=P^T$ caracterizan completamente a las
proyecciones ortogonales.
\end{proof}

\begin{entrega}{Proyecciones oblicuas: por qué se pide $P=P^T$}
\label{ent:oblicuas}
El Teorema~\ref{teo:proyector} caracteriza las proyecciones ortogonales con
\emph{dos} condiciones, $P^2=P$ y $P=P^T$. Este ejercicio explora qué queda si
se conserva sólo la primera.

\begin{enumerate}[label=\textbf{(\alph*)},itemsep=3pt,topsep=4pt]
\item Sea $P\in\R^{n\times n}$ con $P^2=P$ (\emph{idempotente}), sin suponer
simetría. Probar que $\R^n=C(P)\oplus N(P)$ y que, en esa descomposición,
$P\bx$ es la componente de $\bx$ en $C(P)$ \emph{a lo largo de} $N(P)$.
Probar también que $I-P$ es idempotente, con $C(I-P)=N(P)$ y $N(I-P)=C(P)$.
La suma es directa, pero en general \emph{no} es ortogonal.

\item Verificarlo en
$P=\begin{bsmallmatrix}1&1\\ 0&0\end{bsmallmatrix}$: comprobar $P^2=P$,
identificar $C(P)$ y $N(P)$, dibujarlos, y marcar en el dibujo cómo se
proyecta un vector cualquiera.

\item Para $P_t=\begin{bsmallmatrix}1&t\\ 0&0\end{bsmallmatrix}$ con
$t\in\R$, calcular $\|P_t\bx\|$ para $\bx=(0,1)^T$ y mostrar que el cociente
$\|P_t\bx\|/\|\bx\|$ se puede hacer tan grande como se quiera. Comparar con
el caso ortogonal, donde siempre $\|P\bx\|\le\|\bx\|$ (¿por qué?). \emph{Esta
es la razón de fondo para exigir $P=P^T$: las proyecciones oblicuas
amplifican, y todo lo que amplifica amplifica también los errores.}

\item Probar que, para $P$ idempotente, $P=P^T$ si y sólo si
$N(P)=C(P)^{\perp}$.

\item Construir una matriz idempotente de $2\times2$ cuyo espacio columna y
espacio nulo formen un ángulo de $30^{\circ}$, y calcular para ella el máximo
de $\|P\bx\|/\|\bx\|$. ¿Qué pasa con ese máximo cuando el ángulo tiende a
cero?
\end{enumerate}

\medskip
\noindent{\footnotesize\emph{Se entrega por escrito. El ítem (c) explica por
qué en todo el curso ``proyección'' significa ``proyección ortogonal'' salvo
aviso en contrario.}}
\end{entrega}

\begin{observacion}[La proyección complementaria]\label{obs:I-P}
Toda proyección parte a $\bb$ en dos piezas perpendiculares:
\[
\bb=\underbrace{P\bb}_{\text{en } C(Q)}
+\underbrace{(I-P)\bb}_{\text{ortogonal a } C(Q)} .
\]
La matriz $I-P$ \emph{también} es una proyección ortogonal, y proyecta sobre
el complemento ortogonal $C(Q)^{\perp}$. En efecto, cumple las dos
condiciones del Teorema~\ref{teo:proyector}:
\[
(I-P)^T=I-P^T=I-P ,
\qquad
(I-P)^2=I-2P+P^2=I-2P+P=I-P ,
\]
usando $P^T=P$ y $P^2=P$. Y su espacio columna es $C(Q)^{\perp}$, porque
$(I-P)\bb$ es el error, que ya vimos que es ortogonal a $C(Q)$; recíprocamente,
si $\bi{w}\perp C(Q)$ entonces $P\bi{w}=\bi{0}$ y por lo tanto
$\bi{w}=(I-P)\bi{w}\in C(I-P)$.

Las dos proyecciones son, entonces, las dos mitades de la identidad:
$P+(I-P)=I$. Esto ya había aparecido en la Sección~\ref{sec:cuatro-ort} al
descomponer un vector en su parte del espacio fila y su parte del espacio
nulo.
\end{observacion}

\begin{ejemplo}\label{ej:proyecciones}
Tres matrices con columnas ortonormales, creciendo de $3\times1$ a
$3\times2$ y a $3\times3$:
\[
Q_1=\frac13\begin{bmatrix}2\\2\\-1\end{bmatrix}\qquad
Q_2=\frac13\begin{bmatrix}2&2\\2&-1\\-1&2\end{bmatrix}\qquad
Q_3=\frac13\begin{bmatrix}2&2&-1\\2&-1&2\\-1&2&2\end{bmatrix}.
\]
Proyectemos el mismo $\bb=(3,3,3)^T$ sobre los tres espacios columna.

\emph{Sobre la recta} $C(Q_1)$:
\[
P_1\bb=Q_1(Q_1^T\bb)
=\frac13\begin{bmatrix}2\\2\\-1\end{bmatrix}
\cdot\frac{2\cdot3+2\cdot3-1\cdot3}{3}
=\begin{bmatrix}2\\2\\-1\end{bmatrix},
\qquad
\bi{e}=\begin{bmatrix}1\\1\\4\end{bmatrix},
\quad \|\bi{e}\|=\sqrt{18}.
\]

\emph{Sobre el plano} $C(Q_2)$ (que contiene a la recta anterior):
\[
P_2\bb=Q_2(Q_2^T\bb)=\begin{bmatrix}4\\1\\1\end{bmatrix},
\qquad
\bi{e}=\begin{bmatrix}-1\\2\\2\end{bmatrix},
\quad \|\bi{e}\|=\sqrt{9}=3 .
\]
El error se achicó ($3<\sqrt{18}$): el plano contiene a la recta, así que su
punto más cercano a $\bb$ está más cerca.

\emph{Sobre todo el espacio} $C(Q_3)=\R^3$: ahora $Q_3$ es cuadrada,
$P_3=Q_3Q_3^T=I$, y $P_3\bb=\bb$. \textbf{El error es cero.} Proyectar
sobre el espacio total no cambia nada. \hfill$\square$
\end{ejemplo}

\subsection*{¿Y si la base no es ortonormal?}

Todo lo anterior supone que tenemos a mano una base \emph{ortonormal} del
subespacio. En la práctica lo habitual es que el subespacio venga dado como
$C(A)$, con las columnas de $A$ independientes pero no ortonormales. La
fórmula de la proyección sigue existiendo y es
\begin{equation}\label{eq:proyeccion-general}
\boxed{\ P=A(A^T A)^{-1}A^T\ }
\end{equation}
que tiene sentido gracias a la Observación~\ref{obs:ATA} ($A^TA$ es
invertible exactamente cuando las columnas de $A$ son independientes) y que
se reduce a $QQ^T$ cuando $A=Q$ ya tiene columnas ortonormales, porque
entonces $Q^TQ=I$. Demostrarla, y entender por qué es la proyección correcta,
es el siguiente ejercicio de desarrollo.

\begin{entrega}{La proyección sobre $C(A)$ con una base cualquiera}
\label{ent:proyeccion}
Sea $A\in\R^{m\times n}$ con columnas linealmente independientes y sea
$P=A(A^TA)^{-1}A^T$.

\begin{enumerate}[label=\textbf{(\alph*)},itemsep=3pt,topsep=4pt]
\item Probar que $A^TA$ es invertible, de modo que $P$ está bien definida.
\emph{Sugerencia:} es de $n\times n$; alcanza con ver que $N(A^TA)=N(A)$.

\item Probar que $P^2=P$ y $P^T=P$.

\item Probar que $C(P)=C(A)$ y que $N(P)=C(A)^{\perp}$. Deducir que
$\rg(P)=n$.

\item Probar que, para todo $\bb\in\R^m$, el error $\bb-P\bb$ es ortogonal a
$C(A)$, y que $P\bb$ es el punto de $C(A)$ más cercano a $\bb$:
\[
\|\bb-P\bb\| \le \|\bb-\bv\| \qquad\text{para todo } \bv\in C(A),
\]
con igualdad sólo si $\bv=P\bb$. \emph{Sugerencia:} Pitágoras, como en el
Teorema~\ref{teo:proyector}.

\item \emph{(No depende de la base.)} Sea $M\in\R^{n\times n}$ invertible y
$B=AM$ (las columnas de $B$ son otra base del mismo subespacio). Probar que
$B(B^TB)^{-1}B^T = A(A^TA)^{-1}A^T$. Es decir: la proyección depende del
\emph{subespacio}, no de la base elegida para describirlo.

\item Verificar que si las columnas de $A$ son ortonormales, entonces
\eqref{eq:proyeccion-general} se reduce a $P=AA^T$.

\item Probar que $\tr(P)=n$. \emph{Sugerencia:} usar
$\tr(XY)=\tr(YX)$ del Capítulo~\ref{cap:algebra-lineal}. (Es un caso
particular del Corolario~\ref{cor:traza-proyeccion}: para cualquier
proyección, la traza es el rango.)

\item Calcular $P$ para
\[
A=\begin{bmatrix} 1&0\\ 1&1\\ 0&1\end{bmatrix}
\]
y verificar sobre esa matriz concreta los ítems (b), (c) y (g).
\end{enumerate}

\medskip
\noindent{\footnotesize\emph{Se entrega por escrito. La fórmula
\eqref{eq:proyeccion-general} es la que resuelve el problema de mínimos
cuadrados, que veremos en su propio capítulo; el ítem (e) explica por qué
tiene sentido hablar de ``la'' proyección sobre un subespacio.}}
\end{entrega}
\subsection*{Volver a $A\bx=\bb$}

Con las proyecciones en la mano podemos cerrar la discusión que quedó abierta
en el Capítulo~\ref{cap:LU}: allí vimos que el conjunto solución de
$A\bx=\bb$, cuando no es vacío, es $\bx_p+N(A)$, un trasladado del espacio
nulo. La ortogonalidad permite decir dos cosas más: cómo reconocer sin
resolver si el sistema es compatible, y cuál es, entre todas las soluciones,
la \emph{mejor}.

\begin{entrega}{Compatibilidad y la solución de norma mínima}
\label{ent:norma-minima}
Sea $A\in\R^{m\times n}$ y $\bb\in\R^m$.

\begin{enumerate}[label=\textbf{(\alph*)},itemsep=3pt,topsep=4pt]
\item \emph{(Alternativa de Fredholm.)} Probar que
\[
A\bx=\bb \ \text{ tiene solución}
\qquad\Longleftrightarrow\qquad
\by^T\bb=0 \ \text{ para todo } \by \text{ con } A^T\by=\bi{0}.
\]
\emph{Sugerencia:} $C(A)=N(A^T)^{\perp}$, por el Teorema~\ref{teo:fila-nulo}
y la Proposición~\ref{prop:descomposicion}(iv).

\item Interpretarlo en términos de ecuaciones: cada $\by\in N(A^T)$ da una
combinación lineal de las filas de $A$ cuyo resultado es la ecuación
$0=\by^T\bb$; el sistema es compatible exactamente cuando todas esas
``ecuaciones vacías'' son ciertas. Encontrar todas las condiciones de
compatibilidad del sistema
\[
\begin{bmatrix} 1&2\\ 2&4\\ 1&3\end{bmatrix}\bx = \bb ,
\]
calculando primero $N(A^T)$.

\item Supongamos ahora el sistema compatible. Probar que todas sus soluciones
tienen la \emph{misma} componente en el espacio fila: si $\bx$ y $\bx'$ son
soluciones y las escribimos como
$\bx=\bx_{\text{fila}}+\bx_{\text{nulo}}$ según el
Corolario~\ref{coro:cuatro}, entonces
$\bx_{\text{fila}}=\bx'_{\text{fila}}$.

\item Deducir que hay exactamente \emph{una} solución que pertenece a
$C(A^T)$, y que es la de \textbf{norma mínima}: para toda solución $\bx$,
\[
\|\bx\|^2=\|\bx_{\text{fila}}\|^2+\|\bx_{\text{nulo}}\|^2
\ \ge\ \|\bx_{\text{fila}}\|^2 ,
\]
con igualdad si y sólo si $\bx=\bx_{\text{fila}}$.

\item Probar que la restricción de $A$ al espacio fila,
\[
A\big|_{C(A^T)}:\ C(A^T)\longrightarrow C(A),
\]
es una biyección. \emph{(Esta es la semilla de la inversa generalizada y de
las bases especiales $A\bv_k=\sigma_k\bu_k$ de la SVD.)}

\item Para
\[
A=\begin{bmatrix}1&1&1\end{bmatrix},\qquad \bb=\begin{bmatrix}3\end{bmatrix},
\]
describir todas las soluciones y hallar la de norma mínima de dos maneras:
proyectando una solución cualquiera sobre $C(A^T)$, y buscándola
directamente de la forma $\bx=A^T\bi{z}$. Comprobar que dan lo mismo y
explicar por qué toda solución de la forma $A^T\bi{z}$ está en el espacio
fila.
\end{enumerate}

\medskip
\noindent{\footnotesize\emph{Se entrega por escrito. El ítem (a) convierte la
condición $\bb\in C(A)$, que es difícil de verificar, en una ortogonalidad
que se chequea; los ítems (d) y (e) son el punto de partida de la inversa
generalizada.}}
\end{entrega}

% ===================================================================
\section{Ortogonalización de Gram-Schmidt con el producto interno usual en $\R^n$}\label{sec:gram-schmidt}

Aquí hay un hecho clave: \textbf{todo subespacio de $\R^n$ tiene una base
ortogonal} (y, normalizando, una ortonormal). Pensemos en un plano de $\R^3$
generado por dos vectores independientes $\ba$ y $\bb$. Para conseguir una
base ortogonal del mismo plano, a $\bb$ le restamos su componente en la
dirección de $\ba$ (¡una proyección!):
\begin{equation}\label{eq:gs-idea}
\text{Base ortogonal:}\quad \ba
\qquad\text{y}\qquad
\bi{c}=\bb-\frac{\ba^T\bb}{\ba^T\ba}\,\ba .
\end{equation}
La cuenta lo confirma:
$\ba^T\bi{c}=\ba^T\bb-\dfrac{\ba^T\bb}{\ba^T\ba}\,
\ba^T\ba=0$. Y $\bi{c}\ne\bi{0}$ porque $\bb$ no era múltiplo de $\ba$.
El plano generado no cambió: solo cambiamos la base.

Esta idea de ``ortogonalizar'' se aplica a cualquier cantidad de vectores,
restando en cada paso las proyecciones sobre los vectores ya construidos.
Ese es el \textbf{proceso de Gram-Schmidt}.

\begin{teorema}[Gram-Schmidt]\label{teo:gram-schmidt}
Sean $\ba_1,\dots,\ba_r$ linealmente independientes. Definimos
\[
\bi{v}_1=\ba_1,\qquad
\bi{v}_k=\ba_k-\sum_{j=1}^{k-1}
\frac{\bi{v}_j^T\ba_k}{\bi{v}_j^T\bi{v}_j}\,\bi{v}_j
\qquad (k=2,\dots,r).
\]
Entonces $\bi{v}_1,\dots,\bi{v}_r$ son ortogonales, no nulos, y
\[
\hbox{span}\{\bi{v}_1,\dots,\bi{v}_k\}=\hbox{span}\{\ba_1,\dots,\ba_k\}
\qquad\text{para cada }k .
\]
Normalizando, $\bq_k=\bi{v}_k/\|\bi{v}_k\|$ es una base ortonormal del
subespacio generado. En particular, \emph{todo subespacio tiene una base
ortonormal}.
\end{teorema}

\begin{proof}
Por inducción en $k$. Para $k=1$ no hay nada que probar. Supongamos que
$\bi{v}_1,\dots,\bi{v}_{k-1}$ ya son ortogonales, no nulos, y generan lo
mismo que $\ba_1,\dots,\ba_{k-1}$. El vector $\bi{v}_k$ es $\ba_k$ menos su
proyección sobre el subespacio generado por
$\bi{v}_1,\dots,\bi{v}_{k-1}$ (proyectar sobre cada eje ortogonal y sumar,
como en la Sección~\ref{sec:proyecciones}). Para cada $i<k$:
\[
\bi{v}_i^T\bi{v}_k
=\bi{v}_i^T\ba_k-\sum_{j=1}^{k-1}
\frac{\bi{v}_j^T\ba_k}{\bi{v}_j^T\bi{v}_j}\,
\bi{v}_i^T\bi{v}_j
=\bi{v}_i^T\ba_k
-\frac{\bi{v}_i^T\ba_k}{\bi{v}_i^T\bi{v}_i}\,
\bi{v}_i^T\bi{v}_i=0 ,
\]
porque en la suma solo sobrevive el término $j=i$. Además
$\bi{v}_k\ne\bi{0}$: si fuera cero, $\ba_k$ sería combinación de
$\bi{v}_1,\dots,\bi{v}_{k-1}$, es decir de $\ba_1,\dots,\ba_{k-1}$,
contradiciendo la independencia. Los generados coinciden porque $\bi{v}_k$
difiere de $\ba_k$ en una combinación de los anteriores.
\end{proof}

En la práctica conviene normalizar a medida que se avanza: así cada paso es
``restar proyecciones sobre los $\bq$ ya construidos y normalizar''.

\begin{caja}
\textbf{Algoritmo de Gram-Schmidt} (entrada: $\ba_1,\dots,\ba_r$
independientes; salida: $\bq_1,\dots,\bq_r$ ortonormales).
\begin{enumerate}[leftmargin=*,itemsep=1pt]
\item Inicializar: $\bq_1=\ba_1/\|\ba_1\|$.
\item Para $k=2,\dots,r$: restar de $\ba_k$ su proyección sobre lo ya
construido,
\[
\bi{v}_k=\ba_k-\bq_1(\bq_1^T\ba_k)-\cdots
-\bq_{k-1}(\bq_{k-1}^T\ba_k),
\]
y normalizar: $\bq_k=\bi{v}_k/\|\bi{v}_k\|$.
\end{enumerate}
En cada paso, $\hbox{span}\{\bq_1,\dots,\bq_k\}=\hbox{span}\{\ba_1,\dots,\ba_k\}$: el
subespacio construido crece de a una dimensión por vez.
\end{caja}

\begin{ejemplo}[Gram-Schmidt en acción]\label{ej:gs}
Ortogonalicemos
\[
\ba_1=\begin{bmatrix}1\\-1\\0\end{bmatrix},\qquad
\ba_2=\begin{bmatrix}2\\0\\-2\end{bmatrix},\qquad
\ba_3=\begin{bmatrix}3\\-3\\3\end{bmatrix}.
\]
\emph{Paso 1.} $\bi{v}_1=\ba_1=(1,-1,0)^T$.

\emph{Paso 2.} $\bi{v}_1^T\ba_2=2$ y $\bi{v}_1^T\bi{v}_1=2$:
\[
\bi{v}_2=\ba_2-\frac{2}{2}\,\bi{v}_1
=\begin{bmatrix}2\\0\\-2\end{bmatrix}-\begin{bmatrix}1\\-1\\0\end{bmatrix}
=\begin{bmatrix}1\\1\\-2\end{bmatrix}.
\]

\emph{Paso 3.} $\bi{v}_1^T\ba_3=6$, $\bi{v}_2^T\ba_3=-6$,
$\bi{v}_2^T\bi{v}_2=6$:
\[
\bi{v}_3=\ba_3-\frac{6}{2}\,\bi{v}_1-\frac{-6}{6}\,\bi{v}_2
=\begin{bmatrix}3\\-3\\3\end{bmatrix}
-\begin{bmatrix}3\\-3\\0\end{bmatrix}
+\begin{bmatrix}1\\1\\-2\end{bmatrix}
=\begin{bmatrix}1\\1\\1\end{bmatrix}.
\]
Verificación: $\bi{v}_1^T\bi{v}_2=1-1+0=0$,
$\bi{v}_1^T\bi{v}_3=0$, $\bi{v}_2^T\bi{v}_3=0$. Normalizando,
\[
\bq_1=\frac{1}{\sqrt2}\begin{bmatrix}1\\-1\\0\end{bmatrix},\qquad
\bq_2=\frac{1}{\sqrt6}\begin{bmatrix}1\\1\\-2\end{bmatrix},\qquad
\bq_3=\frac{1}{\sqrt3}\begin{bmatrix}1\\1\\1\end{bmatrix}
\]
es una base ortonormal de $\R^3$. \hfill$\square$
\end{ejemplo}

\begin{observacion}
En la Sección~\ref{sec:QR} veremos que este proceso, escrito en forma
matricial, \emph{es} una factorización: si las columnas de $A$ son
$\ba_1,\dots,\ba_r$ y las de $Q$ son $\bq_1,\dots,\bq_r$, entonces
$A=QR$ con $R$ triangular superior.
\end{observacion}

% ===================================================================
\section{Matrices  reales ortogonales}\label{sec:matrices-ortogonales}

\subsection*{Altas y delgadas primero: $Q^T Q=I$}

Sea $Q$ de $m\times r$ con columnas ortonormales $\bq_1,\dots,\bq_r$, es
decir,
\[
\bq_i^T\bq_j=\delta_{ij}=
\begin{cases} 1 & \text{si } i=j,\\ 0 & \text{si } i\neq j.\end{cases}
\]
Esas $r^2$ condiciones se resumen en una sola igualdad matricial. En efecto,
las filas de $Q^T$ son las columnas de $Q$ traspuestas, de modo que al
multiplicar $Q^T$ por $Q$ la entrada $(i,j)$ del producto es
\[
(Q^TQ)_{ij}=\fila_i(Q^T)\,\col_j(Q)=\bq_i^T\bq_j=\delta_{ij} ,
\]
es decir, \emph{$Q^TQ$ es la tabla de todos los productos escalares entre
columnas de $Q$}. Pedir que las columnas sean ortonormales es pedir que esa
tabla sea la identidad:
\[
Q^T Q=
\begin{bmatrix}\text{---}\ \bq_1^T\ \text{---}\\ \vdots\\
\text{---}\ \bq_r^T\ \text{---}\end{bmatrix}
\begin{bmatrix}\bq_1\ \cdots\ \bq_r\end{bmatrix}
=\begin{bmatrix}
\bq_1^T\bq_1 & \cdots & \bq_1^T\bq_r\\
\vdots & \ddots & \vdots\\
\bq_r^T\bq_1 & \cdots & \bq_r^T\bq_r
\end{bmatrix}
=\begin{bmatrix}1&&\\&\ddots&\\&&1\end{bmatrix}=I_r .
\]
Conviene leer las dos lecturas por separado: los \emph{unos} de la diagonal
dicen que cada columna es unitaria, y los \emph{ceros} de afuera dicen que
cada par de columnas distintas es ortogonal. En resumen,
\[
\boxed{\ \text{las columnas de } Q \text{ son ortonormales}
\iff Q^TQ=I_r\ }
\]

Vale la pena insistir en que esto \emph{no} dice que $Q$ sea invertible. La
igualdad $Q^TQ=I_r$ dice que $Q^T$ es inversa \emph{a izquierda} de $Q$, y
nada más; cuando $m>r$ la matriz $Q$ ni siquiera es cuadrada. En ese caso se
la llama \textbf{semiortogonal}, y el producto en el otro orden
\emph{no} da la identidad: $QQ^T\neq I_m$. De hecho ya sabemos qué es, porque
es la proyección $P$ del Teorema~\ref{teo:proyector}, una matriz de rango $r$;
si fuera $QQ^T=I_m$ tendría rango $m>r$, imposible. Dicho de otro modo: las
$m$ filas de $Q$ son $m$ vectores de $\R^r$ y no pueden ser ortonormales si
$m>r$, porque en $\R^r$ no entran más de $r$ vectores ortonormales.

\begin{proposicion}[$Q$ preserva normas y ángulos]\label{prop:preserva}
Si $Q^T Q=I$ entonces, para todos $\bx,\by$,
\[
\|Q\bx\|=\|\bx\|
\qquad\text{y}\qquad
(Q\bx)^T(Q\by)=\bx^T\by .
\]
\end{proposicion}

\begin{proof}
$(Q\bx)^T(Q\by)=\bx^T Q^T Q\,\by=\bx^T\by$; con
$\by=\bx$ queda $\|Q\bx\|^2=\|\bx\|^2$.
\end{proof}

Multiplicar por $Q$ no cambia magnitudes ni ángulos. Esta es la razón por la
que las matrices ortogonales son las favoritas del cálculo numérico:
\textbf{los cálculos con $Q$ nunca se desbordan}. Multiplicar muchas veces
por matrices ortogonales no agranda ni achica los errores de redondeo.

\subsection*{Matrices ortogonales: cuadradas}

\begin{definicion}
Una \textbf{matriz ortogonal} es una matriz \emph{cuadrada} con columnas
ortonormales.
\end{definicion}

\begin{teorema}\label{teo:Q-inversa}
Si $Q$ es cuadrada, son equivalentes:
\begin{enumerate}[label=(\roman*),itemsep=1pt]
\item las columnas de $Q$ son ortonormales, es decir $Q^T Q=I$;
\item $Q$ es invertible y $Q^{-1}=Q^T$;
\item $QQ^T=I$, es decir las filas de $Q$ son ortonormales;
\item $\|Q\bx\|=\|\bx\|$ para todo $\bx$.
\end{enumerate}
\end{teorema}

\begin{proof}
(i)$\Rightarrow$(ii): $Q$ cuadrada con columnas ortonormales tiene columnas
independientes (Proposición~\ref{prop:ort-indep}), luego es invertible; de
$Q^T Q=I$, la inversa a izquierda $Q^T$ es \emph{la} inversa.
(ii)$\Rightarrow$(iii): $QQ^T=QQ^{-1}=I$, y ese producto calcula los
productos internos entre filas: las filas son ortonormales (¡probablemente
una base ortonormal distinta de la de las columnas!).
(iii)$\Rightarrow$(i): aplicando (i)$\Rightarrow$(iii) a $Q^T$.
(i)$\Rightarrow$(iv) es la Proposición~\ref{prop:preserva}. Para
(iv)$\Rightarrow$(i): de $\|Q\bx\|^2=\|\bx\|^2$ para todo $\bx$ se
recuperan los productos internos por polarización,
$\bx^T\by=\tfrac12(\|\bx+\by\|^2-\|\bx\|^2-\|\by\|^2)$, aplicada a
$Q\bx$ y $Q\by$; en particular $\bq_i^T\bq_j=\bi{e}_i^T\bi{e}_j$.
\end{proof}

\begin{observacion}
El nombre honesto sería ``matriz ortonormal'', pero la tradición manda:
cuando $Q$ es cuadrada y $Q^T=Q^{-1}$, se dice matriz ortogonal.
\end{observacion}

\begin{observacion}[La identidad como suma de matrices de rango uno]
\label{obs:identidad-rango-uno}
Escribamos $QQ^T=I$ con la multiplicación
\emph{columnas por filas} del Capítulo~\ref{cap:algebra-lineal}:
\[
I=\bq_1\bq_1^T+\bq_2\bq_2^T+\cdots+\bq_n\bq_n^T .
\]
¡La identidad descompuesta en $n$ matrices de rango uno, una por cada eje
de la base ortonormal! El caso trivial es
$I=\sum_k\bi{e}_k\bi{e}_k^T$ con la base canónica. Aplicada a un
vector, esta identidad es otra vez el Teorema~\ref{teo:coeficientes}:
\[
\bv=I\bv=\bq_1(\bq_1^T\bv)+\cdots+\bq_n(\bq_n^T\bv)
=c_1\bq_1+\cdots+c_n\bq_n .
\]
Y si en lugar de sumar los $n$ términos sumamos solo $r$, obtenemos
$\bq_1\bq_1^T+\cdots+\bq_r\bq_r^T=QQ^T$ con $Q$ de
$n\times r$: la proyección de la Sección~\ref{sec:proyecciones}.
\emph{Proyectar es truncar la identidad.}
\end{observacion}

\begin{proposicion}[Productos e inversas]\label{prop:productos}
Si $Q_1$ y $Q_2$ son ortogonales del mismo tamaño, entonces $Q_1Q_2$ es
ortogonal, y $Q^{-1}=Q^T$ es ortogonal.
\end{proposicion}

\begin{proof}
$(Q_1Q_2)^T(Q_1Q_2)=Q_2^T Q_1^T Q_1Q_2=Q_2^T Q_2=I$.
Y $(Q^T)^T Q^T=QQ^T=I$ por el
Teorema~\ref{teo:Q-inversa}.
\end{proof}

\subsection*{Matrices ortogonales famosas}

\textbf{1. Rotaciones.} Cuando todo el plano gira un ángulo $\theta$
alrededor de $(0,0)$, las magnitudes no cambian y los ángulos entre vectores
tampoco. Las columnas son las imágenes de $\bi{e}_1$ y $\bi{e}_2$: vectores
unitarios ortogonales, con $\cos^2\theta+\sin^2\theta=1$:
\begin{equation}\label{eq:rotacion}
Q_{\text{rota}}=\begin{bmatrix}\cos\theta&-\sin\theta\\
\sin\theta&\cos\theta\end{bmatrix}
=\textbf{rotación de ángulo } \theta,
\qquad \det Q_{\text{rota}}=+1 .
\end{equation}

\textbf{2. Reflexiones.} Si multiplicamos la segunda columna por $-1$, las
dos columnas siguen siendo ortogonales y unitarias:
\begin{equation}\label{eq:reflexion}
Q_{\text{refleja}}=\begin{bmatrix}\cos\theta&\sin\theta\\
\sin\theta&-\cos\theta\end{bmatrix}
=\textbf{reflexión en la recta a ángulo } \tfrac{\theta}{2},
\qquad \det Q_{\text{refleja}}=-1 .
\end{equation}
Esta $Q$ refleja cada vector como en un espejo: la recta espejo forma ángulo
$\theta/2$ con el eje $x$. La rotación tiene determinante $+1$ (el plano
gira); la reflexión tiene determinante $-1$ (el plano se da vuelta).

\begin{figure}[ht]
\centering
\begin{tikzpicture}[scale=1.5,>=stealth]
% rotacion
\begin{scope}
\draw[->,gray] (-1.3,0)--(1.4,0);
\draw[->,gray] (0,-0.3)--(0,1.4);
\draw[->,thick] (0,0)--(1,0) node[below right]{\small $\bi{e}_1$};
\draw[->,thick] (0,0)--(0,1) node[above left]{\small $\bi{e}_2$};
\draw[->,thick,blue] (0,0)--(40:1)
  node[right]{\small $Q\bi{e}_1=\binom{\cos\theta}{\sin\theta}$};
\draw[->,thick,blue] (0,0)--(130:1)
  node[above left]{\small $Q\bi{e}_2=\binom{-\sin\theta}{\cos\theta}$};
\draw[->] (0.55,0) arc (0:40:0.55); \node at (18:0.75) {\small $\theta$};
\draw[->] (90:0.4) arc (90:130:0.4);
\node at (0,-0.75) {\small Rotación: $\det=+1$};
\end{scope}
% reflexion
\begin{scope}[shift={(4.6,0)}]
\draw[->,gray] (-1.3,0)--(1.4,0);
\draw[->,gray] (0,-1.1)--(0,1.4);
\draw[dashed] (-1.1,0) ++(20:2.2) -- (200:1.1)
  node[pos=0.05,above]{\small recta a $\theta/2$};
\draw[->,thick] (0,0)--(1,0) node[below right]{\small $\bi{e}_1$};
\draw[->,thick] (0,0)--(0,1) node[above left]{\small $\bi{e}_2$};
\draw[->,thick,blue] (0,0)--(40:1)
  node[right]{\small $Q\bi{e}_1=\binom{\cos\theta}{\sin\theta}$};
\draw[->,thick,blue] (0,0)--(-50:1)
  node[below right]{\small $Q\bi{e}_2=\binom{\sin\theta}{-\cos\theta}$};
\node at (0,-1.45) {\small Reflexión: $\det=-1$};
\end{scope}
\end{tikzpicture}
\caption{Girar todo el plano un ángulo $\theta$; reflejar cada vector
respecto de la recta a ángulo $\theta/2$ (con $\theta=40^{\circ}$). Las
columnas de $Q$ son $Q\bi{e}_1$ y $Q\bi{e}_2$.}
\label{fig:rota-refleja}
\end{figure}

Componer transformaciones es multiplicar matrices, y la
Proposición~\ref{prop:productos} dice que el producto sigue siendo
ortogonal: rotación por rotación $=$ rotación; reflexión por reflexión $=$
rotación; rotación por reflexión $=$ reflexión (los determinantes se
multiplican: $(+1)(+1)=+1$, $(-1)(-1)=+1$, $(+1)(-1)=-1$). Todo esto sigue
valiendo en $\R^n$.

\textbf{3. Reflexiones de Householder.} Ejemplos hermosos de reflexiones en
$\R^n$. Se elige un vector unitario $\bu$ y se define
\begin{equation}\label{eq:householder}
H=I-2\,\bu\bu^T .
\end{equation}
$H$ es simétrica (es la identidad menos dos veces la matriz simétrica de
rango uno $\bu\bu^T$) y es ortogonal: usando $\bu^T\bu=1$,
\[
H^T H=H^2=(I-2\bu\bu^T)(I-2\bu\bu^T)
=I-4\bu\bu^T+4\bu\,(\bu^T\bu)\,\bu^T=I .
\]
Dos reflexiones seguidas devuelven cada vector a su lugar. Geométricamente,
$H$ refleja respecto del hiperplano perpendicular a $\bu$:
\[
H\bu=\bu-2\bu(\bu^T\bu)=-\bu ,
\qquad
H\bw=\bw \quad\text{si } \bw\perp\bu .
\]
Los ``autovalores'' de $H$ son $-1$ (una vez, en la dirección $\bu$) y $+1$
($n-1$ veces, en el espejo): como toda reflexión.

\begin{observacion}[Reflexiones $\leftrightarrow$ proyecciones]
\label{obs:reflexion-proyeccion}
La construcción de Householder es un caso particular de una
correspondencia general. Si $P$ es \emph{cualquier} proyección ortogonal
($P^2=P=P^T$), la matriz
\[
H=I-2P
\]
es simétrica y ortogonal:
$H^2=(I-2P)^2=I-4P+4P^2=I$. Y refleja: para $\bw\in C(P)$ es
$H\bw=\bw-2\bw=-\bw$, mientras que para $\bi{z}\perp C(P)$ es
$H\bi{z}=\bi{z}$. Es la reflexión que invierte $C(P)$ y deja fijo su
complemento ortogonal. Householder es el caso $P=\bu\bu^T$
(proyección sobre la recta de $\bu$). Recíprocamente, si $H$ es simétrica
con $H^2=I$, entonces
\[
P=\tfrac12(I-H)
\]
cumple $P^T=P$ y
$P^2=\tfrac14(I-2H+H^2)=\tfrac14(2I-2H)=P$: es una proyección ortogonal.
\textbf{Reflexiones y proyecciones son dos caras de la misma moneda}:
$H=I-2P$, $P=\tfrac12(I-H)$.
\end{observacion}

Con $\bu=(1,1,\dots,1)^T/\sqrt n$, la fórmula
$H_n=I-\frac{2}{n}\,\textbf{unos}(n,n)$ da ejemplos fáciles de recordar:
\[
H_3=\frac13\begin{bmatrix}1&-2&-2\\-2&1&-2\\-2&-2&1\end{bmatrix}
\qquad
H_4=\frac12\begin{bmatrix}1&-1&-1&-1\\-1&1&-1&-1\\-1&-1&1&-1\\
-1&-1&-1&1\end{bmatrix},
\]
y $H_4$ se parece a una matriz de Hadamard.

\begin{entrega}{Reflexiones de Householder}\label{ent:householder}
Sea $\bu\in\R^n$ con $\|\bu\|=1$ y sea $H=I-2\,\bu\bu^T$.

\begin{enumerate}[label=\textbf{(\alph*)},itemsep=3pt,topsep=4pt]
\item Probar que $H^T=H$, que $H^2=I$ y que $H$ es ortogonal, y concluir que
$H^{-1}=H$: \emph{reflejar dos veces es no hacer nada}.

\item Probar que $H\bu=-\bu$ y que $H\bw=\bw$ para todo $\bw\perp\bu$.
Describir en palabras qué le hace $H$ a un vector cualquiera, descomponiéndolo
según $\R^n=\gen\{\bu\}\oplus\bu^{\perp}$.

\item \emph{(El lema que hace funcionar todo.)} Sean $\bx\neq\by$ con
$\|\bx\|=\|\by\|$. Probar que, tomando
\[
\bu=\frac{\bx-\by}{\|\bx-\by\|},
\]
resulta $H\bx=\by$. \emph{Sugerencia:} calcular $\|\bx-\by\|^2$ usando
$\|\bx\|=\|\by\|$ y ver que $\bu^T\bx=\tfrac12\|\bx-\by\|$.

\item Aplicar el ítem anterior al caso que nos interesa: dado
$\bx\neq\bi{0}$, los únicos
$\by$ múltiplos de $\bi{e}_1$ con $\|\by\|=\|\bx\|$ son
$\by=\pm\|\bx\|\,\bi{e}_1$. Escribir $\bu$ explícitamente en cada caso.

\item Hacer las cuentas con $\bx=(1,2,2)^T$ para las dos elecciones de signo
y verificar que
\begin{align*}
H&=\frac13\begin{bmatrix}1&2&2\\2&1&-2\\2&-2&1\end{bmatrix}
&&\text{da}& H\bx&=3\,\bi{e}_1,\\[2pt]
H'&=\frac13\begin{bmatrix}-1&-2&-2\\-2&2&-1\\-2&-1&2\end{bmatrix}
&&\text{da}& H'\bx&=-3\,\bi{e}_1 .
\end{align*}

\item \emph{(Estabilidad --- opcional.)} En el ítem (c) hay que dividir por
$\|\bx-\by\|$. Mostrar que, con $\by=\pm\|\bx\|\bi{e}_1$,
\[
\|\bx-\by\|^2 = 2\|\bx\|\big(\|\bx\|\mp x_1\big),
\]
y explicar por qué conviene elegir el signo
$\by=-\sgn(x_1)\,\|\bx\|\,\bi{e}_1$. ¿Qué pasa si se elige el otro y $\bx$
ya está casi alineado con $\bi{e}_1$?

\item \emph{(Costo --- opcional.)} Para $X\in\R^{n\times k}$, comparar el costo de
calcular $HX$ ensamblando primero la matriz $H$ con el de calcular
$X-\bu\,(2\bu^TX)$ sin ensamblarla nunca. \emph{Las reflexiones de
Householder no se guardan como matrices: se guarda el vector $\bu$.}

\end{enumerate}

\medskip
\noindent{\footnotesize\emph{Se entrega por escrito. Los ítems (c), (d) y (e)
son exactamente lo que necesita la construcción de la factorización $QR$ de la
Sección~\ref{sec:QR}: allí los usaremos sin volver a demostrarlos. Los ítems
(f) y (g) son opcionales: explican, respectivamente, cómo se elige el signo y
por qué estas reflexiones nunca se guardan como matrices.}}
\end{entrega}

\textbf{4. Permutaciones.} Ya las usamos: son las matrices que aparecieron en
la Sección~\ref{sec:pivoteo} del Capítulo~\ref{cap:LU} para intercambiar filas
durante el pivoteo, y las que hacen que la factorización $PA=LU$ exista
siempre. Una \textbf{matriz de permutación} $P$ tiene las mismas columnas que
la identidad, en algún orden; sus columnas son vectores canónicos distintos y
por lo tanto ortonormales. Allá probamos $P^{-1}=P^T$ a mano, y dijimos que
las matrices con esa propiedad se llamarían ortogonales: acá se cierra esa
promesa, porque es exactamente la definición de arriba. La inversa es también
una permutación: la que deshace el reordenamiento. Por ejemplo,
\[
P=\begin{bmatrix}0&1&0&0\\0&0&1&0\\0&0&0&1\\1&0&0&0\end{bmatrix}
\qquad\text{tiene columnas ortonormales, }P^T P=I .
\]

\textbf{5. Rotaciones de Givens.} Son rotaciones que actúan sobre un solo par
de coordenadas y dejan todas las demás quietas. Dados $i<j$ y un ángulo
$\theta$, la matriz $G(i,j,\theta)\in\R^{n\times n}$ es la identidad salvo
las cuatro entradas
\[
g_{ii}=g_{jj}=c=\cos\theta,
\qquad
g_{ij}=-s=-\sin\theta,
\qquad
g_{ji}=s=\sin\theta,
\]
es decir, es la rotación \eqref{eq:rotacion} embebida en las filas y columnas
$i,j$:
\[
G(2,4,\theta)=
\begin{bmatrix}
1&0&0&0\\
0&c&0&-s\\
0&0&1&0\\
0&s&0&c
\end{bmatrix}.
\]
Multiplicar $G\bx$ sólo cambia las coordenadas $i$ y $j$ de $\bx$, rotándolas
en su plano. Eligiendo bien el ángulo se puede \emph{anular} una de las dos:
mientras Householder pone en cero toda una columna de un solo golpe, Givens
lo hace entrada por entrada. Eso lo vuelve la herramienta adecuada cuando la
matriz ya tiene muchos ceros y no queremos destruirlos.

\begin{entrega}{Rotaciones de Givens}\label{ent:givens}
\begin{enumerate}[label=\textbf{(\alph*)},itemsep=3pt,topsep=4pt]
\item Probar que $G(i,j,\theta)$ es ortogonal y que $\det G=+1$.
Describir $G^T$ en términos de $G$.

\item Dado $\bx\in\R^n$ con $(x_i,x_j)\neq(0,0)$, hallar $c$ y $s$ con
$c^2+s^2=1$ tales que $G(i,j,\theta)\bx$ tenga un cero en la coordenada $j$,
y comprobar que en la coordenada $i$ queda $\sqrt{x_i^2+x_j^2}$. Escribir
$c$ y $s$ sin usar funciones trigonométricas, y explicar cómo calcular
$\sqrt{x_i^2+x_j^2}$ sin que se desborde cuando $x_i$ o $x_j$ son muy
grandes. \emph{Sugerencia:} sacar factor común el mayor de los dos en
valor absoluto.

\item Verificar que las demás coordenadas de $\bx$ no cambian, y que
$\|G\bx\|=\|\bx\|$ (dos veces: por la Proposición~\ref{prop:preserva} y
haciendo la cuenta).

\item \emph{(Factorización $QR$ por Givens.)} Describir cómo anular, una por
una, todas las entradas debajo de la diagonal de una matriz
$A\in\R^{m\times n}$ usando rotaciones de Givens, y verificar que se obtiene
$G_N\cdots G_1A=\widetilde{R}$ y por lo tanto
$A=\widetilde{Q}\widetilde{R}$ con
$\widetilde{Q}=G_1^T\cdots G_N^T$. ¿Cuántas rotaciones hacen falta si $A$ es
llena? Hacerlo explícitamente con
\[
A=\begin{bmatrix} 3&1\\ 4&1\\ 0&2\end{bmatrix}.
\]

\item \emph{(Cuándo conviene.)} Estimar el costo de $QR$ por Givens para una
matriz llena y compararlo con Householder. Después, contar cuántas rotaciones
hacen falta si $A$ es \emph{tridiagonal} de $n\times n$, o si es de Hessenberg
superior (ceros debajo de la primera subdiagonal). Explicar por qué en esos
casos Givens gana por lejos.

\item \emph{(Actualizar una factorización.)} Supongamos conocida
$A=\widetilde{Q}\widetilde{R}$ y que se agrega una fila nueva $\bi{a}^T$ al
final de $A$. Mostrar cómo recuperar la factorización de la matriz ampliada
con $n$ rotaciones de Givens, en lugar de refactorizar desde cero. \emph{Este
es el procedimiento que se usa cuando los datos llegan de a una observación
por vez.}
\end{enumerate}

\medskip
\noindent{\footnotesize\emph{Se entrega por escrito. Las rotaciones de Givens
vuelven a aparecer en el algoritmo $QR$ para calcular autovalores, donde las
matrices que hay que factorizar son justamente de Hessenberg.}}
\end{entrega}

\begin{observacion}[De dónde salen los vectores ortogonales]
Cuando una matriz es simétrica u ortogonal, \textbf{tiene autovectores
ortogonales}. Esta es la fuente más importante de vectores ortogonales de la
matemática aplicada, y es el corazón de los Capítulos~\ref{cap:schur} a
\ref{cap:SVD}: teorema espectral $S=Q\Lambda Q^T$ y SVD $A=U\Sigma V^T$. Los ejercicios
sobre la matriz de Fourier (autovectores de una permutación) son el primer
ejemplo de este fenómeno.
\end{observacion}

% ===================================================================
\section{La factorización $A=QR$}\label{sec:QR}

Cerramos el capítulo juntando las dos ideas centrales: ortonormalizar
(Gram-Schmidt) y reflejar (Householder). Las dos producen la misma
factorización.

\subsection*{Gram-Schmidt en forma matricial}

Sea $A\in\R^{m\times n}$ con columnas $\ba_1,\dots,\ba_n$ linealmente
independientes, y apliquémosles Gram-Schmidt para obtener
$\bq_1,\dots,\bq_n$. La clave es leer al revés el algoritmo: en el paso $k$
escribimos
\[
\ba_k=(\bq_1^T\ba_k)\,\bq_1+\cdots+(\bq_{k-1}^T\ba_k)\,\bq_{k-1}
+\|\bi{v}_k\|\,\bq_k ,
\]
porque $\bi{v}_k=\ba_k-\sum_{j<k}(\bq_j^T\ba_k)\bq_j$ y
$\bq_k=\bi{v}_k/\|\bi{v}_k\|$. Es decir: \emph{cada $\ba_k$ es combinación
lineal de los primeros $k$ vectores $\bq$}, y los coeficientes son productos
internos.

Acá es donde conviene recordar la fórmula \eqref{eq:cols} del
Capítulo~\ref{cap:algebra-lineal}, que decía
\[
\col_k(QR)=Q\,\col_k(R) .
\]
Leída de derecha a izquierda, dice exactamente lo que necesitamos: \emph{una
combinación lineal de las columnas de $Q$ es una columna de un producto
$QR$, y los coeficientes de esa combinación son la columna correspondiente de
$R$}. Nuestra relación de arriba es precisamente una combinación lineal de
$\bq_1,\dots,\bq_k$ que da $\ba_k$, de modo que si llamamos $\bi{r}_k$ al
vector de coeficientes,
\[
\bi{r}_k=
\big(\bq_1^T\ba_k,\ \dots,\ \bq_{k-1}^T\ba_k,\ \|\bi{v}_k\|,\ 0,\dots,0\big)^T
\in\R^n ,
\]
entonces $\ba_k=Q\,\bi{r}_k$ para cada $k$. Poniendo esas $n$ igualdades una
al lado de la otra, es decir armando la matriz cuya columna $k$ es
$\bi{r}_k$, obtenemos de un saque
\[
\underbrace{[\,\ba_1\ \cdots\ \ba_n\,]}_{A}
=Q\,\underbrace{[\,\bi{r}_1\ \cdots\ \bi{r}_n\,]}_{R} ,
\qquad\text{o sea}\qquad
A=QR,
\]
donde los ceros al final de cada $\bi{r}_k$ son los que vuelven triangular
superior a $R$. Explícitamente:
\[
Q=\begin{bmatrix}\bq_1&\cdots&\bq_n\end{bmatrix}\in\R^{m\times n},
\qquad
R=\begin{bmatrix}
\|\bi{v}_1\| & \bq_1^T\ba_2 & \bq_1^T\ba_3 & \cdots & \bq_1^T\ba_n\\
0 & \|\bi{v}_2\| & \bq_2^T\ba_3 & \cdots & \bq_2^T\ba_n\\
0 & 0 & \|\bi{v}_3\| & \cdots & \bq_3^T\ba_n\\
\vdots & & & \ddots & \vdots\\
0 & 0 & \cdots & 0 & \|\bi{v}_n\|
\end{bmatrix}\in\R^{n\times n}.
\]
La matriz $R$ es triangular superior \emph{porque} $\ba_k$ no usa
$\bq_{k+1},\dots,\bq_n$: el subespacio generado crece de a una dimensión por
vez. Y su diagonal es positiva porque $\|\bi{v}_k\|>0$.

\begin{teorema}[Factorización $QR$ reducida]\label{teo:QR}
Sea $A\in\R^{m\times n}$ con columnas linealmente independientes (en
particular $m\ge n$). Entonces existen $Q\in\R^{m\times n}$ con columnas
ortonormales y $R\in\R^{n\times n}$ triangular superior con diagonal
positiva tales que
\[
A=QR .
\]
Además $Q$ y $R$ son \emph{únicas} con esas condiciones.
\end{teorema}

\begin{proof}
La existencia es la construcción de arriba. Para la unicidad, supongamos
$Q_1R_1=Q_2R_2$ con las dos $R_i$ triangulares superiores de diagonal
positiva (y por lo tanto invertibles). Entonces
$Q_2^TQ_1=R_2R_1^{-1}$. El miembro izquierdo es ortogonal: en efecto,
$Q_1=Q_2(R_2R_1^{-1})$ y multiplicando por su traspuesta,
$I=Q_1^TQ_1=(R_2R_1^{-1})^T(R_2R_1^{-1})$. Llamemos $T=R_2R_1^{-1}$: es
triangular superior con diagonal positiva (producto e inversa de
triangulares superiores de diagonal positiva) y cumple $T^TT=I$. Pero
entonces $T^{-1}=T^T$ es a la vez triangular superior y triangular inferior,
o sea diagonal; y una matriz diagonal con $T^TT=I$ y diagonal positiva es
$T=I$. Luego $R_2=R_1$ y $Q_2=Q_1$.
\end{proof}

\begin{ejemplo}[$QR$ del Ejemplo~\ref{ej:gs}]
Con $\ba_1=(1,-1,0)^T$, $\ba_2=(2,0,-2)^T$, $\ba_3=(3,-3,3)^T$ habíamos
obtenido $\bi{v}_1=(1,-1,0)^T$, $\bi{v}_2=(1,1,-2)^T$,
$\bi{v}_3=(1,1,1)^T$, de normas $\sqrt2,\sqrt6,\sqrt3$. Los coeficientes de
$R$ son $\bq_1^T\ba_2=\sqrt2$, $\bq_1^T\ba_3=3\sqrt2$ y
$\bq_2^T\ba_3=-\sqrt6$, así que
\[
\underbrace{\begin{bmatrix}1&2&3\\-1&0&-3\\0&-2&3\end{bmatrix}}_{A}
=
\underbrace{\begin{bmatrix}
\tfrac{1}{\sqrt2}&\tfrac{1}{\sqrt6}&\tfrac{1}{\sqrt3}\\[2pt]
-\tfrac{1}{\sqrt2}&\tfrac{1}{\sqrt6}&\tfrac{1}{\sqrt3}\\[2pt]
0&-\tfrac{2}{\sqrt6}&\tfrac{1}{\sqrt3}
\end{bmatrix}}_{Q}
\underbrace{\begin{bmatrix}
\sqrt2&\sqrt2&3\sqrt2\\ 0&\sqrt6&-\sqrt6\\ 0&0&\sqrt3
\end{bmatrix}}_{R} .
\]
Conviene verificar columna por columna: la primera dice
$\ba_1=\sqrt2\,\bq_1$; la segunda, $\ba_2=\sqrt2\,\bq_1+\sqrt6\,\bq_2$.
\hfill$\square$
\end{ejemplo}

\begin{observacion}[$QR$ completa]\label{obs:QR-completa}
La $Q$ del Teorema~\ref{teo:QR} es \emph{alta y delgada}: tiene columnas
ortonormales pero no es una matriz ortogonal si $m>n$. Completando
$\bq_1,\dots,\bq_n$ hasta una base ortonormal de $\R^m$ se obtiene la
\textbf{factorización $QR$ completa}
\[
A=\widetilde{Q}\,\widetilde{R},
\qquad
\widetilde{Q}\in\R^{m\times m}\ \text{ortogonal},
\qquad
\widetilde{R}=\begin{bmatrix}R\\ 0\end{bmatrix}\in\R^{m\times n},
\]
donde las $m-n$ columnas nuevas de $\widetilde{Q}$ son una base ortonormal
de $C(A)^{\perp}=N(A^T)$. Las dos versiones son útiles: la reducida es la que
se guarda, la completa es la que se usa en las demostraciones.
\end{observacion}

\subsection*{$QR$ por reflexiones de Householder}

Gram-Schmidt construye $Q$ ortonormalizando las columnas de $A$. La
construcción de Householder hace exactamente lo contrario: deja $A$ quieta y
la va \emph{triangulando} a fuerza de reflexiones, igual que la eliminación
gaussiana la triangula a fuerza de restas de filas --- con la diferencia
capital de que ahora las operaciones son ortogonales y por lo tanto no
amplifican los errores.

La herramienta son los ítems (c) y (d) del Ejercicio de
desarrollo~\ref{ent:householder}: dado $\bx$, existe una reflexión $H$ con
\[
H\bx=-\sgn(x_1)\,\|\bx\|\,\bi{e}_1 ,
\]
es decir, una reflexión que concentra todo $\bx$ en su primera coordenada y
anula las demás. Apliquémosla a la primera columna de $A$:
\[
H_1A=\begin{bmatrix}
\rho_1 & \n{r}_1^T\\
\bi{0} & A_1
\end{bmatrix},
\qquad \rho_1=-\sgn(a_{11})\|\col_1(A)\| ,
\]
con $A_1\in\R^{(m-1)\times(n-1)}$. Ahora repetimos sobre $A_1$: tomamos
$H_2'\in\R^{(m-1)\times(m-1)}$ que anula todo debajo del primer elemento de
$\col_1(A_1)$ y la embebemos sin tocar la primera fila ni la primera columna,
\[
H_2=\begin{bmatrix}1&\bi{0}^T\\ \bi{0}&H_2'\end{bmatrix},
\]
que sigue siendo ortogonal y simétrica. Después de $n$ pasos,
\[
H_n\cdots H_2H_1\,A=\begin{bmatrix}R\\ 0\end{bmatrix},
\qquad\text{y por lo tanto}\qquad
A=\underbrace{H_1H_2\cdots H_n}_{\widetilde{Q}}
\begin{bmatrix}R\\ 0\end{bmatrix},
\]
usando que cada $H_k$ es su propia inversa. El producto de matrices
ortogonales es ortogonal (Proposición~\ref{prop:productos}), así que
$\widetilde{Q}$ lo es.

\begin{teorema}[$QR$ por Householder]\label{teo:QR-householder}
Toda matriz $A\in\R^{m\times n}$ con $m\ge n$ admite una factorización
$A=\widetilde{Q}\widetilde{R}$ con $\widetilde{Q}\in\R^{m\times m}$ ortogonal
y $\widetilde{R}=\begin{bsmallmatrix}R\\ 0\end{bsmallmatrix}$ con $R$
triangular superior de $n\times n$. Si las columnas de $A$ son
independientes, $R$ es invertible.
\end{teorema}

\begin{observacion}[Cuál conviene]
Las dos construcciones dan la misma factorización cuando las columnas son
independientes (por la unicidad del Teorema~\ref{teo:QR}, salvo signos de las
columnas de $Q$), pero no son equivalentes en la práctica.
\begin{itemize}[itemsep=2pt]
\item Gram-Schmidt es más transparente y da directamente la versión reducida,
pero es \emph{numéricamente frágil}: los $\bq_k$ calculados pierden
ortogonalidad a medida que $k$ crece, y el efecto empeora cuanto más
parecidas entre sí son las columnas de $A$. (La variante llamada
\emph{Gram-Schmidt modificado}, que resta las proyecciones de a una en lugar
de todas juntas, es matemáticamente idéntica y numéricamente mucho mejor.)
\item Householder trabaja siempre con matrices ortogonales, que por la
Proposición~\ref{prop:preserva} no agrandan nada, y produce una $\widetilde{Q}$
ortogonal hasta la precisión de la máquina. Además no hay que ensamblar
ninguna $H_k$: alcanza con guardar los vectores $\bu_k$.
\item No hace falta que las columnas de $A$ sean independientes: Householder
funciona igual y los ceros aparecen solos en la diagonal de $R$.
\end{itemize}
Por eso todas las bibliotecas numéricas calculan $QR$ con reflexiones de
Householder.
\end{observacion}

\begin{observacion}[Para qué sirve]
La factorización $QR$ es la herramienta estándar para resolver el problema de
\textbf{mínimos cuadrados}: minimizar $\|A\bx-\bb\|$. Como multiplicar por
$\widetilde{Q}^T$ no cambia las normas,
\[
\|A\bx-\bb\|^2=\|\widetilde{Q}^T(A\bx-\bb)\|^2
=\Big\|\begin{bmatrix}R\\ 0\end{bmatrix}\bx
-\begin{bmatrix}\bi{c}_1\\ \bi{c}_2\end{bmatrix}\Big\|^2
=\|R\bx-\bi{c}_1\|^2+\|\bi{c}_2\|^2 ,
\]
con $\bi{c}=\widetilde{Q}^T\bb$. El primer término se puede hacer cero
resolviendo el sistema triangular $R\bx=\bi{c}_1$, y el segundo es el error
mínimo, que no depende de $\bx$. Desarrollamos esto en el capítulo dedicado a
mínimos cuadrados, donde lo compararemos con las ecuaciones normales
$A^TA\bx=A^T\bb$ que salen de la proyección
\eqref{eq:proyeccion-general}. Allí se verá también por qué esta vía es
preferible: formar $A^TA$ eleva al cuadrado el número de condición del
problema.
\end{observacion}

% ===================================================================
\section*{Ideas clave en $\R^n$}
\begin{ideas}
\begin{itemize}[leftmargin=*]
\item \textbf{Ortogonalidad}: $\bx\perp\by\iff\bx^T\by=0$ (caso
complejo: $\overline{\bx}^T\by=0$). Pitágoras:
$\|\bx\pm\by\|^2=\|\bx\|^2+\|\by\|^2$. Cauchy-Schwarz:
$|\bx^T\by|\le\|\bx\|\,\|\by\|$, que define el ángulo (similitud
coseno) $\cos\theta=\bx^T\by/(\|\bx\|\,\|\by\|)$.
\item \textbf{Base ortonormal} $\bq_1,\dots,\bq_n$: los coeficientes de
$\bv=\sum c_k\bq_k$ se obtienen por separado, $c_k=\bq_k^T\bv$
(matricialmente, $\bv=Q\bi{c}\Rightarrow\bi{c}=Q^T\bv$).
\item \textbf{Subespacios ortogonales}: $C(A^T)\perp N(A)$ y
$C(A)\perp N(A^T)$; las dimensiones suman $n$ y $m$. Todo
$\bv=\bv_{\text{fila}}+\bv_{\text{nulo}}$, con componentes ortogonales; en
general $\R^n=\mathbf{W}\oplus\mathbf{W}^{\perp}$ y
$(\mathbf{W}^{\perp})^{\perp}=\mathbf{W}$ (ortogonales $\ne$
complementos: las dimensiones deben sumar $n$).
\item \textbf{Proyección}: con base ortonormal del subespacio,
$P=QQ^T$ cumple $P=P^T=P^2$; el error $(I-P)\bb$ es ortogonal al
subespacio y $P\bb$ es el punto más cercano. Recíprocamente,
$P^2=P=P^T$ \emph{caracteriza} a las proyecciones ortogonales, y
$H=I-2P$ da la reflexión asociada. La identidad se descompone como
$I=\sum_k\bq_k\bq_k^T$: proyectar es truncar esa suma.
\item \textbf{Gram-Schmidt}: restar proyecciones ortogonaliza. Todo
subespacio tiene base ortonormal.
\item \textbf{Matriz ortogonal}: cuadrada con $Q^T Q=I$, o sea
$Q^{-1}=Q^T$; preserva normas y ángulos ($\|Q\bx\|=\|\bx\|$). Famosas:
rotaciones ($\det=+1$), reflexiones y Householder $H=I-2\bu\bu^T$
($\det=-1$), rotaciones de Givens, permutaciones, Hadamard$/\sqrt n$,
Fourier (compleja).
Productos e inversas de ortogonales son ortogonales.
\item \textbf{Factorización $A=QR$}: $Q$ con columnas ortonormales, $R$
triangular superior de diagonal positiva; es Gram-Schmidt leído como
producto de matrices, y es única. La misma factorización se obtiene
triangulando $A$ con reflexiones de Householder, que es la forma
numéricamente estable y la que usan las bibliotecas.
\item \textbf{Anticipo}: simétricas y ortogonales tienen autovectores
ortogonales; proyecciones y $QR$ $\to$ mínimos cuadrados.
\end{itemize}
\end{ideas}

% ===================================================================
\section*{Ejercicios}

\noindent\footnotesize Las etiquetas indican el sabor de cada ejercicio:
\textsc{[Cálculo]} para fijar la mecánica, \textsc{[Concepto]} para
justificar y demostrar, \textsc{[Datos]} para la conexión con ciencia de
datos. Las soluciones están al final del capítulo.\normalsize

\begin{enumerate}[leftmargin=*,label=\textbf{\thechapter.\arabic*}]

\item \tipo{Cálculo} Si $\bu$ y $\bv$ son vectores unitarios ortogonales,
demuestre que $\bu+\bv$ es ortogonal a $\bu-\bv$. ¿Cuáles son las magnitudes
de esos vectores?

\item \tipo{Cálculo} Dibuje dos vectores unitarios $\bu$ y $\bv$ que
\emph{no sean} ortogonales. Muestre que $\bw=\bv-\bu(\bu^T\bv)$ es
ortogonal a $\bu$ (y agregue $\bw$ al gráfico). ¿Qué paso de Gram-Schmidt
acaba de hacer?

\item \tipo{Concepto} Dibuje dos vectores $\bu$ y $\bv$ cualesquiera a
partir del origen $(0,0)$. Complete dos lados más para hacer un
paralelogramo con diagonales $\bw=\bu+\bv$ y $\bi{z}=\bu-\bv$. Muestre que
$\bw^T\bw+\bi{z}^T\bi{z}$ es igual a $2\bu^T\bu+2\bv^T\bv$
(\emph{ley del paralelogramo}: la suma de los cuadrados de las diagonales es
la suma de los cuadrados de los cuatro lados).

\item \tipo{Cálculo} Sean $\bx=(1,i)^T$ y $\by=(i,1)^T$ en
$\Cx^2$.
\begin{enumerate}[label=(\alph*)]
\item Calcule $\bx^T\by$ y $\overline{\bx}^T\by$. ¿Son $\bx$ y
$\by$ ortogonales?
\item Calcule $\|\bx\|$ y $\|\by\|$ con la norma correcta
$\|\bx\|^2=\overline{\bx}^T\bx$. ¿Qué daría la fórmula real
$\bx^T\bx$, y por qué es inaceptable?
\end{enumerate}

\item \tipo{Concepto} Suponga que $H_n$ es una matriz de Hadamard: entradas
$\pm1$ y columnas ortogonales, con $n\ge2$.
\begin{enumerate}[label=(\alph*)]
\item Mirando dos columnas, muestre que $n$ debe ser par.
\item Mirando tres columnas, muestre que $n$ debe ser múltiplo de $4$.
\emph{Ayuda:} puede suponer que la primera columna es toda de unos
(¿por qué?); cuente con $a$ = cantidad de filas donde la segunda y la
tercera columna tienen ambas un $+1$, etc.
\end{enumerate}

\item\label{ej:proy-recta} \tipo{Cálculo} Sea $\bq=\frac13(2,2,-1)^T$ (unitario) y
$P=\bq\bq^T$.
\begin{enumerate}[label=(\alph*)]
\item Escriba la matriz $P$ y verifique $P=P^T=P^2$. ¿Cuál es su rango?
\item Proyecte $\bb=(0,3,3)^T$ sobre la recta generada por $\bq$ y
verifique que el error $\bi{e}=\bb-P\bb$ es ortogonal a $\bq$.
\item ¿Sobre qué subespacio proyecta $I-P$? Verifique
$\|P\bb\|^2+\|\bi{e}\|^2=\|\bb\|^2$.
\end{enumerate}

\item \tipo{Cálculo} Aplique Gram-Schmidt a
\[
\ba_1=\begin{bmatrix}1\\0\\1\end{bmatrix},\qquad
\ba_2=\begin{bmatrix}1\\1\\0\end{bmatrix},\qquad
\ba_3=\begin{bmatrix}0\\1\\1\end{bmatrix},
\]
y escriba la base ortonormal resultante $\bq_1,\bq_2,\bq_3$.

\item \tipo{Concepto} Propiedad clave de toda matriz ortogonal:
$\|Q\bx\|^2=\|\bx\|^2$ para todo vector $\bx$. Más aún, muestre que
$(Q\bx)^T(Q\by)=\bx^T\by$ para todo par $\bx,\by$. Entonces $Q$
no modifica las magnitudes ni los ángulos. ¡Los cálculos con $Q$ nunca se
desbordan! Explique por qué esto también garantiza que las columnas de $Q$
(las imágenes $Q\bi{e}_k$) son ortonormales.

\item \tipo{Cálculo} Con las matrices de \eqref{eq:rotacion} y
\eqref{eq:reflexion}:
\begin{enumerate}[label=(\alph*)]
\item Multiplique dos rotaciones de ángulos $\theta_1$ y $\theta_2$ y
muestre que el producto es la rotación de ángulo $\theta_1+\theta_2$
(recupere las fórmulas de suma de senos y cosenos).
\item Sin multiplicar matrices, use determinantes para decidir: ¿reflexión
por reflexión es rotación o reflexión? ¿Y rotación por reflexión?
\end{enumerate}

\item \tipo{Concepto} En $\R^2$, tome el vector unitario
$\bu=(-\sin\varphi,\cos\varphi)^T$, normal a la recta que forma ángulo
$\varphi$ con el eje $x$. Calcule $H=I-2\bu\bu^T$ y muestre que
coincide con la reflexión \eqref{eq:reflexion} con $\theta=2\varphi$: las
reflexiones de Householder en $\R^2$ son exactamente las matrices
$Q_{\text{refleja}}$.

\item \tipo{Cálculo} Una \textbf{matriz de permutación} tiene las mismas
columnas que la matriz identidad (en algún orden). \emph{Explique por qué
esta y todas las matrices de permutación son ortogonales}:
\[
P=\begin{bmatrix}0&1&0&0\\0&0&1&0\\0&0&0&1\\1&0&0&0\end{bmatrix}
\quad\text{tiene columnas ortonormales, entonces } P^T P=I
\text{ y } P^{-1}=P^T .
\]
¿Qué permutación es $P^{-1}$, en términos de la que hace $P$?

\item \tipo{Cálculo} Cuatro autovectores de esa matriz $P$ son
$\bx_1=(1,1,1,1)$, $\bx_2=(1,i,i^2,i^3)$, $\bx_3=(1,i^2,i^4,i^6)$ y
$\bx_4=(1,i^3,i^6,i^9)$. Multiplique $P$ por cada uno de esos vectores para
encontrar $\lambda_1,\lambda_2,\lambda_3,\lambda_4$. Los autovectores son
las columnas de la \textbf{matriz de Fourier $\bi{F}$} de $4\times4$.
\[
\text{Muestre que } Q=\frac{F}{2}=\frac12
\begin{bmatrix}1&1&1&1\\1&i&-1&i^3\\1&i^2&1&i^6\\1&i^3&-1&i^9\end{bmatrix}
\text{ tiene columnas ortonormales}:\ \overline{Q}^T Q=I .
\]
(La matriz de Fourier es la matriz ortogonal compleja más importante de la
matemática aplicada. Como es compleja, el test correcto es
$\overline{\bx}_j^T\bx_k=0$.)

\item \tipo{Cálculo} \textbf{Las wavelets de Haar son vectores ortogonales
(columnas de $W$) que usan solo $1,-1,0$}:
\[
n=4\qquad
W=\begin{bmatrix}1&1&1&0\\1&1&-1&0\\1&-1&0&1\\1&-1&0&-1\end{bmatrix}
\qquad
\begin{array}{l}
\textbf{Encuentre } W^T W,\ W^{-1}\textbf{ y las}\\
\textbf{ocho wavelets de Haar para } n=8 .
\end{array}
\]
Las wavelets de Haar pueden pensarse como funciones con soporte en el
intervalo $[0,1]$. Son estudiadas en diferentes campos de la matemática y
resultan una herramienta útil en informática e ingeniería, por ejemplo para
la codificación de señales.

\begin{center}
\begin{tabular}{c@{\hspace{18pt}}c@{\hspace{18pt}}c@{\hspace{18pt}}c}
\begin{tikzpicture}[scale=1.1]
\draw[step=0.5,gray!30,very thin] (-0.4,-1.4) grid (1.4,1.4);
\draw[->,thick] (-0.45,0)--(1.45,0);
\draw[->,thick] (0,-1.45)--(0,1.45);
\draw[gray,thick] (0,0)--(0,1);
\draw[blue,very thick] (0,1)--(1,1);
\draw[gray,thick] (1,1)--(1,0);
\end{tikzpicture}
&
\begin{tikzpicture}[scale=1.1]
\draw[step=0.5,gray!30,very thin] (-0.4,-1.4) grid (1.4,1.4);
\draw[->,thick] (-0.45,0)--(1.45,0);
\draw[->,thick] (0,-1.45)--(0,1.45);
\draw[gray,thick] (0,0)--(0,1);
\draw[blue,very thick] (0,1)--(0.5,1);
\draw[gray,thick] (0.5,1)--(0.5,-1);
\draw[blue,very thick] (0.5,-1)--(1,-1);
\draw[gray,thick] (1,-1)--(1,0);
\end{tikzpicture}
&
\begin{tikzpicture}[scale=1.1]
\draw[step=0.5,gray!30,very thin] (-0.4,-1.4) grid (1.4,1.4);
\draw[->,thick] (-0.45,0)--(1.45,0);
\draw[->,thick] (0,-1.45)--(0,1.45);
\draw[gray,thick] (0,0)--(0,1);
\draw[blue,very thick] (0,1)--(0.25,1);
\draw[gray,thick] (0.25,1)--(0.25,-1);
\draw[blue,very thick] (0.25,-1)--(0.5,-1);
\draw[gray,thick] (0.5,-1)--(0.5,0);
\draw[blue,very thick] (0.5,0)--(1,0);
\end{tikzpicture}
&
\begin{tikzpicture}[scale=1.1]
\draw[step=0.5,gray!30,very thin] (-0.4,-1.4) grid (1.4,1.4);
\draw[->,thick] (-0.45,0)--(1.45,0);
\draw[->,thick] (0,-1.45)--(0,1.45);
\draw[blue,very thick] (0,0)--(0.5,0);
\draw[gray,thick] (0.5,0)--(0.5,1);
\draw[blue,very thick] (0.5,1)--(0.75,1);
\draw[gray,thick] (0.75,1)--(0.75,-1);
\draw[blue,very thick] (0.75,-1)--(1,-1);
\draw[gray,thick] (1,-1)--(1,0);
\end{tikzpicture}
\end{tabular}
\end{center}

\item \tipo{Datos} (Compresión con la base de Haar.) Normalice las columnas
de la matriz $W$ del ejercicio anterior para obtener una base ortonormal
$\bw_1,\bw_2,\bw_3,\bw_4$ de $\R^4$, y tome la ``señal''
$\bv=(5,3,-1,1)^T$.
\begin{enumerate}[label=(\alph*)]
\item Calcule los cuatro coeficientes $c_k=\bw_k^T\bv$ (¡cada uno por
separado!) y verifique la reconstrucción $\bv=\sum c_k\bw_k$.
\item Verifique la identidad de Parseval $\|\bv\|^2=c_1^2+c_2^2+c_3^2+c_4^2$
y explique por qué vale en cualquier base ortonormal.
\item \emph{Comprima}: conserve solo los dos coeficientes más grandes y
reconstruya $\widetilde{\bv}$. ¿Qué forma tiene $\widetilde{\bv}$ y cuánto
vale el error $\|\bv-\widetilde{\bv}\|$? (Esta es la idea de la codificación
de señales e imágenes: guardar pocos coeficientes en una base ortonormal
bien elegida.)
\end{enumerate}

\item \tipo{Datos} Sea $Q$ ortogonal de $n\times n$ y sean
$\bx_1,\dots,\bx_N\in\R^n$ los puntos de un conjunto de datos.
\begin{enumerate}[label=(\alph*)]
\item Muestre que $\|Q\bx_i-Q\bx_j\|=\|\bx_i-\bx_j\|$ para todo par $i,j$:
rotar (o reflejar) la nube de datos no cambia ninguna distancia.
\item Concluya que cualquier método basado en distancias (por ejemplo,
buscar el vecino más cercano de un punto) da exactamente el mismo resultado
con los datos $\bx_i$ que con los datos rotados $Q\bx_i$.
\item Si en cambio multiplicamos por una matriz invertible \emph{no}
ortogonal, dé un ejemplo de $2\times2$ donde el vecino más cercano de un
punto cambie.
\end{enumerate}

\item \tipo{Datos} (Similitud coseno.) Tres ``documentos'' se representan
por sus conteos de tres palabras clave:
$\bx=(3,1,0)^T$, $\by=(6,2,0)^T$, $\bi{z}=(0,2,3)^T$.
\begin{enumerate}[label=(\alph*)]
\item Calcule las tres similitudes coseno entre pares.
\item ¿Por qué $\bx$ e $\by$ tienen similitud exactamente $1$? Relacione
con el caso de igualdad de Cauchy-Schwarz.
\item Explique por qué la similitud coseno es la medida natural cuando un
documento puede ser ``el doble de largo'' que otro con el mismo tema.
\end{enumerate}

\item \tipo{Concepto} (Reflexiones $\leftrightarrow$ proyecciones.) Sea
$P$ una proyección ortogonal ($P^2=P=P^T$) y $H=I-2P$.
\begin{enumerate}[label=(\alph*)]
\item Muestre que $H$ es simétrica y ortogonal, que $H\bw=-\bw$ para
$\bw\in C(P)$ y que $H\bi{z}=\bi{z}$ para $\bi{z}\perp C(P)$.
\item Recíprocamente, si $H$ es simétrica con $H^2=I$, muestre que
$P=\frac12(I-H)$ es una proyección ortogonal.
\item Tome $P=\bq\bq^T$ con $\bq=\frac13(2,2,-1)^T$ (ejercicio
\textbf{5.6}) y escriba la matriz $H=I-2P$. ¿Qué hace $H$
geométricamente?
\end{enumerate}

\end{enumerate}

% ===================================================================
\section*{Soluciones}

\begin{enumerate}[leftmargin=*,label=\textbf{\thechapter.\arabic*}]

\item $(\bu+\bv)^T(\bu-\bv)=\bu^T\bu-\bu^T\bv+\bv^T\bu
-\bv^T\bv=\|\bu\|^2-\|\bv\|^2=1-1=0$. Las magnitudes:
$\|\bu\pm\bv\|^2=\|\bu\|^2+\|\bv\|^2=2$ (Pitágoras), así que ambas
diagonales miden $\sqrt2$.

\item $\bu^T\bw=\bu^T\bv-(\bu^T\bu)(\bu^T\bv)
=\bu^T\bv-\bu^T\bv=0$, usando $\bu^T\bu=1$. En el dibujo,
$\bw$ es $\bv$ menos su ``sombra'' sobre $\bu$. Es exactamente el paso 2 de
Gram-Schmidt \eqref{eq:gs-idea}: restar a $\bv$ su proyección sobre $\bu$.

\item
$\bw^T\bw+\bi{z}^T\bi{z}
=(\|\bu\|^2+2\bu^T\bv+\|\bv\|^2)+(\|\bu\|^2-2\bu^T\bv+\|\bv\|^2)
=2\bu^T\bu+2\bv^T\bv$. Los términos cruzados se cancelan entre las
dos diagonales, sean o no ortogonales $\bu$ y $\bv$.

\item (a) $\bx^T\by=1\cdot i+i\cdot1=2i\ne0$, pero el test correcto da
$\overline{\bx}^T\by=1\cdot i+\overline{i}\cdot1=i-i=0$: sí son
ortogonales. (b) $\|\bx\|^2=|1|^2+|i|^2=2$, luego $\|\bx\|=\|\by\|=\sqrt2$.
La fórmula real daría $\bx^T\bx=1+i^2=0$: un vector no nulo con
``norma'' cero, inaceptable; por eso la norma compleja conjuga.

\item (a) Sean $\bi{h}_j,\bi{h}_k$ dos columnas. Cada producto
$h_{ij}h_{ik}$ vale $\pm1$, y la suma de los $n$ productos es $0$: la mitad
son $+1$ y la mitad $-1$, así que $n$ es par. (b) Cambiar de signo una fila
completa no altera ni las entradas $\pm1$ ni la ortogonalidad de las
columnas; haciéndolo fila por fila podemos suponer que la columna 1 es toda
de unos. Entonces, por (a) aplicado contra la columna 1, las columnas 2 y 3
tienen exactamente $n/2$ unos. Sean $a,b,c,d$ las cantidades de filas donde
(columna 2, columna 3) valen $(+,+),(+,-),(-,+),(-,-)$. Ortogonalidad con la
columna 1: $a+b=c+d=n/2$ y $a+c=b+d=n/2$, de donde $b=c$ y $a=d$.
Ortogonalidad entre las columnas 2 y 3: $a-b-c+d=0$, o sea $2a=2b$. Entonces
$a=b=c=d$ y $n=4a$: múltiplo de $4$.

\item (a)
\[
P=\bq\bq^T=\frac19\begin{bmatrix}4&4&-2\\4&4&-2\\-2&-2&1\end{bmatrix},
\]
claramente simétrica, y $P^2=\bq(\bq^T\bq)\bq^T=\bq\bq^T=P$
porque $\bq^T\bq=1$. Es de rango 1 (todas las columnas son múltiplos de
$\bq$). (b) $\bq^T\bb=\frac{0+6-3}{3}=1$, así que
$P\bb=\bq=(2/3,2/3,-1/3)^T$ y
$\bi{e}=\bb-P\bb=(-2/3,7/3,10/3)^T$;
$\bq^T\bi{e}=\frac{-4/3+14/3-10/3}{3}=0$. (c) $I-P$ proyecta sobre el
plano $\bq^{\perp}$, el complemento ortogonal de la recta. Pitágoras:
$\|P\bb\|^2+\|\bi{e}\|^2=1+17=18=\|\bb\|^2$.

\item $\bi{v}_1=(1,0,1)^T$. Después
$\bi{v}_2=\ba_2-\frac{1}{2}\bi{v}_1=(\tfrac12,1,-\tfrac12)^T$
(o, reescalado, $(1,2,-1)^T$). Después, con
$\bi{v}_1^T\ba_3=1$ y $\bi{v}_2^T\ba_3=\tfrac12$,
$\bi{v}_2^T\bi{v}_2=\tfrac32$:
\[
\bi{v}_3=\ba_3-\frac12\,\bi{v}_1-\frac{1/2}{3/2}\,\bi{v}_2
=\left(-\tfrac23,\tfrac23,\tfrac23\right)^T
\quad\text{(o, reescalado, }(-1,1,1)^T\text{)}.
\]
Normalizando:
\[
\bq_1=\frac{1}{\sqrt2}\begin{bmatrix}1\\0\\1\end{bmatrix},\qquad
\bq_2=\frac{1}{\sqrt6}\begin{bmatrix}1\\2\\-1\end{bmatrix},\qquad
\bq_3=\frac{1}{\sqrt3}\begin{bmatrix}-1\\1\\1\end{bmatrix}.
\]

\item $(Q\bx)^T(Q\by)=\bx^T Q^T Q\by=\bx^T\by$; tomando
$\by=\bx$, $\|Q\bx\|^2=\|\bx\|^2$. Como las normas y los productos internos
no cambian, tampoco cambian los ángulos
($\cos\theta=\bx^T\by/\|\bx\|\|\by\|$), y los números nunca crecen al
multiplicar por $Q$. Las columnas de $Q$ son $Q\bi{e}_1,\dots,Q\bi{e}_n$; el
producto interno de dos de ellas es
$(Q\bi{e}_i)^T(Q\bi{e}_j)=\bi{e}_i^T\bi{e}_j$, que vale $1$ si
$i=j$ y $0$ si no: ortonormales.

\item (a) Escribiendo $c_k=\cos\theta_k$ y $s_k=\sin\theta_k$:
\[
\begin{bmatrix}c_1&-s_1\\ s_1&c_1\end{bmatrix}
\begin{bmatrix}c_2&-s_2\\ s_2&c_2\end{bmatrix}
=\begin{bmatrix}c_1c_2-s_1s_2&-(c_1s_2+s_1c_2)\\
s_1c_2+c_1s_2&c_1c_2-s_1s_2\end{bmatrix},
\]
que es la rotación de ángulo $\theta_1+\theta_2$; comparando entradas se
leen las fórmulas de suma
$\cos(\theta_1+\theta_2)=\cos\theta_1\cos\theta_2-\sin\theta_1\sin\theta_2$
y
$\sin(\theta_1+\theta_2)=\sin\theta_1\cos\theta_2+\cos\theta_1\sin\theta_2$. (b) El producto de dos ortogonales es ortogonal,
y su determinante es el producto de los determinantes: reflexión por
reflexión tiene $\det=(-1)(-1)=+1$, es una rotación; rotación por reflexión
tiene $\det=-1$, es una reflexión.

\item Con $\bu=(-\sin\varphi,\cos\varphi)^T$:
\[
\bu\bu^T=\begin{bmatrix}\sin^2\varphi&-\sin\varphi\cos\varphi\\
-\sin\varphi\cos\varphi&\cos^2\varphi\end{bmatrix},
\]
\[
H=I-2\bu\bu^T
=\begin{bmatrix}1-2\sin^2\varphi&2\sin\varphi\cos\varphi\\
2\sin\varphi\cos\varphi&1-2\cos^2\varphi\end{bmatrix}
=\begin{bmatrix}\cos2\varphi&\sin2\varphi\\ \sin2\varphi&-\cos2\varphi
\end{bmatrix},
\]
usando las fórmulas del ángulo doble. Es \eqref{eq:reflexion} con
$\theta=2\varphi$: refleja respecto de la recta a ángulo
$\theta/2=\varphi$, que es justo la recta perpendicular a $\bu$. Las
Householder de $2\times2$ son todas las reflexiones.

\item Las columnas de $P$ son los vectores canónicos
$\bi{e}_4,\bi{e}_1,\bi{e}_2,\bi{e}_3$: distintos vectores de la base
canónica, que es ortonormal. Entonces $P^T P=I$ y $P^{-1}=P^T$,
que vuelve a ser una matriz de permutación: la que \emph{deshace} el
reordenamiento (si $P$ corre cíclicamente las coordenadas hacia arriba,
$P^T$ las corre hacia abajo).

\item Multiplicar por $P$ corre las coordenadas: $P(x_1,x_2,x_3,x_4)^T
=(x_2,x_3,x_4,x_1)^T$. Entonces
\[
P\bx_1=\bx_1\ (\lambda_1=1),\qquad
P\bx_2=(i,i^2,i^3,1)^T=i\,\bx_2\ (\lambda_2=i),
\]
\[
P\bx_3=(i^2,i^4,i^6,1)^T=i^2\bx_3\ (\lambda_3=i^2=-1),\qquad
P\bx_4=i^3\bx_4\ (\lambda_4=i^3=-i).
\]
Los autovalores son las cuatro raíces cuartas de $1$. Para la
ortonormalidad, con $\omega=i$: la entrada $(j,k)$ de
$\overline{F}^T F$ es
\[
\overline{\bx}_j^T\bx_k=\sum_{m=0}^{3}\overline{\omega^{jm}}\,
\omega^{km}=\sum_{m=0}^{3}\omega^{(k-j)m}
=\begin{cases}4&j=k\\[2pt]
\dfrac{\omega^{4(k-j)}-1}{\omega^{\,k-j}-1}=0&j\ne k,\end{cases}
\]
(suma geométrica de raíces de la unidad). Entonces
$\overline{F}^T F=4I$ y $Q=F/2$ cumple $\overline{Q}^T Q=I$.
Por ejemplo,
$\overline{\bx}_1^T\bx_2=1+i+i^2+i^3=(1-1)+(i-i)=0$ y
$\overline{\bx}_2^T\bx_3=(1,-i,-1,i)\,(1,i^2,i^4,i^6)^T
=1+i-1-i=0$: al conjugar el primer vector, $i$ se cambia por $-i$.

\item Las columnas de $W$ son ortogonales (verificarlo par a par), con
normas al cuadrado $4,4,2,2$:
\[
W^T W=\begin{bmatrix}4&0&0&0\\0&4&0&0\\0&0&2&0\\0&0&0&2\end{bmatrix},
\qquad
W^{-1}=(W^T W)^{-1}W^T
=\begin{bmatrix}1/4&1/4&1/4&1/4\\ 1/4&1/4&-1/4&-1/4\\
1/2&-1/2&0&0\\ 0&0&1/2&-1/2\end{bmatrix}.
\]
(Como las columnas son ortogonales pero no unitarias, la inversa es
$D^{-1}W^T$ con $D=W^T W$ diagonal.) Para $n=8$, cada wavelet se
``parte'' en dos mitades de soporte:
\[
W_8=\begin{bmatrix}
1&1&1&0&1&0&0&0\\
1&1&1&0&-1&0&0&0\\
1&1&-1&0&0&1&0&0\\
1&1&-1&0&0&-1&0&0\\
1&-1&0&1&0&0&1&0\\
1&-1&0&1&0&0&-1&0\\
1&-1&0&-1&0&0&0&1\\
1&-1&0&-1&0&0&0&-1
\end{bmatrix}:
\]
la constante; la que separa mitades; dos a escala $1/2$; y cuatro ``finas''
del tipo $(1,-1)$ en posiciones consecutivas.

\item Las columnas normalizadas son
$\bw_1=\tfrac12(1,1,1,1)^T$, $\bw_2=\tfrac12(1,1,-1,-1)^T$,
$\bw_3=\tfrac{1}{\sqrt2}(1,-1,0,0)^T$,
$\bw_4=\tfrac{1}{\sqrt2}(0,0,1,-1)^T$.
(a) $c_1=\tfrac{5+3-1+1}{2}=4$, $c_2=\tfrac{5+3+1-1}{2}=4$,
$c_3=\tfrac{5-3}{\sqrt2}=\sqrt2$, $c_4=\tfrac{-1-1}{\sqrt2}=-\sqrt2$.
Reconstrucción:
$4\bw_1+4\bw_2=(4,4,0,0)^T$, $\sqrt2\,\bw_3=(1,-1,0,0)^T$,
$-\sqrt2\,\bw_4=(0,0,-1,1)^T$; la suma es $(5,3,-1,1)^T=\bv$.
(b) $\|\bv\|^2=25+9+1+1=36$ y $c_1^2+c_2^2+c_3^2+c_4^2=16+16+2+2=36$.
En general $\bv=Q\bi{c}$ con $Q$ ortogonal, y entonces
$\|\bv\|=\|Q\bi{c}\|=\|\bi{c}\|$: las matrices ortogonales preservan
normas. (c) Conservando $c_1=c_2=4$:
$\widetilde{\bv}=(4,4,0,0)^T$, que promedia la señal por bloques
($4=\tfrac{5+3}{2}$, $0=\tfrac{-1+1}{2}$). El error es
$\bv-\widetilde{\bv}=(1,-1,-1,1)^T$ con
$\|\bv-\widetilde{\bv}\|=2=\sqrt{c_3^2+c_4^2}$: por Parseval, el error de
comprimir es exactamente la energía de los coeficientes descartados.

\item (a) $\|Q\bx_i-Q\bx_j\|=\|Q(\bx_i-\bx_j)\|=\|\bx_i-\bx_j\|$ por la
Proposición~\ref{prop:preserva}. (b) Todas las distancias entre pares son
las mismas antes y después, así que cualquier comparación de distancias (en
particular ``¿cuál es el punto más cercano a $\bx_i$?'') da la misma
respuesta. (c) Con $A=\begin{bmatrix}1&0\\0&1/10\end{bmatrix}$ y los puntos
$\bx_1=(0,0)$, $\bx_2=(1,0)$, $\bx_3=(0,2)$: antes, el más cercano a
$\bx_1$ es $\bx_2$ ($1<2$); después, $A\bx_2=(1,0)$ queda a distancia $1$
pero $A\bx_3=(0,0.2)$ queda a distancia $0.2$: el vecino más cercano cambió.
Las matrices invertibles cualesquiera deforman; las ortogonales solo mueven
rígidamente.

\item (a) $\|\bx\|=\sqrt{10}$, $\|\by\|=\sqrt{40}=2\sqrt{10}$,
$\|\bi{z}\|=\sqrt{13}$. Entonces
\[
\cos(\bx,\by)=\frac{18+2}{\sqrt{10}\cdot2\sqrt{10}}=\frac{20}{20}=1,
\qquad
\cos(\bx,\bi{z})=\frac{2}{\sqrt{10}\sqrt{13}}=\frac{2}{\sqrt{130}}
\approx 0{,}18,
\]
y $\cos(\by,\bi{z})=\frac{4}{2\sqrt{10}\sqrt{13}}=\frac{2}{\sqrt{130}}$,
igual que el anterior. (b) $\by=2\bx$: son linealmente dependientes con
constante positiva, que es exactamente el caso de igualdad de
Cauchy-Schwarz (Teorema~\ref{teo:cauchy-schwarz}); mismas proporciones de
palabras $\Rightarrow$ misma dirección $\Rightarrow$ similitud $1$.
(c) Duplicar el largo de un documento duplica su vector de conteos pero no
cambia su dirección: la similitud coseno ignora la magnitud y compara solo
las proporciones (el ``tema''). Por eso $\by$ tiene la misma similitud con
$\bi{z}$ que $\bx$.

\item (a) $H^T=(I-2P)^T=I-2P=H$ y
$H^2=I-4P+4P^2=I-4P+4P=I$, así que $H^T H=H^2=I$: ortogonal. Si
$\bw\in C(P)$ entonces $P\bw=\bw$ (la proyección deja fijo al subespacio)
y $H\bw=\bw-2\bw=-\bw$; si $\bi{z}\perp C(P)$ entonces $P\bi{z}=\bi{0}$ y
$H\bi{z}=\bi{z}$. (b) $P^T=\frac12(I-H)^T=P$ y
\[
P^2=\tfrac14(I-2H+H^2)=\tfrac14(2I-2H)=\tfrac12(I-H)=P :
\]
proyección ortogonal (por la recíproca del
Teorema~\ref{teo:proyector}). (c) Con $P$ del Ejercicio~\ref{ej:proy-recta},
\[
H=I-2P=\frac19\begin{bmatrix}1&-8&4\\-8&1&4\\4&4&7\end{bmatrix},
\]
que es la reflexión de Householder $I-2\bq\bq^T$: invierte la
dirección de $\bq$ y deja fijo el plano espejo $\bq^{\perp}$.

\end{enumerate}

 
